\documentclass[11pt,a4paper]{article}
\usepackage[margin=2.9cm]{geometry}
\usepackage{amsmath,amssymb,amsthm,mathrsfs,bm}
\usepackage[T1]{fontenc}
\usepackage{microtype}
\usepackage{booktabs,longtable,array}
\usepackage[colorlinks=true,linkcolor=blue,citecolor=blue]{hyperref}

% ------------------------------------------------------------------
% Phase V2 of the gate-(II.11) execution (scoping: gate_II11_scoping.tex).
% Contents: A4(a) interaction L^2 comparison on the space-time torus with an
% explicit rate; A4(b) the renormalization dictionary delta_inf = log(4/pi)/(2 pi)
% with rate; A4(c) mean-shift coefficient bounds and convergence; A5 uniform-in-N
% Nelson bounds (Part I Appendix-C pattern).  Fixed throughout: L, m > 0,
% lambda >= 0, and the time-circle circumference t >= t_0 > 0 (constants are
% locally uniform in t and this is tracked).
% ------------------------------------------------------------------

\newcolumntype{P}[1]{>{\raggedright\arraybackslash}p{#1}}

\theoremstyle{plain}
\newtheorem{theorem}{Theorem}[section]
\newtheorem{proposition}[theorem]{Proposition}
\newtheorem{lemma}[theorem]{Lemma}
\newtheorem{corollary}[theorem]{Corollary}
\theoremstyle{definition}
\newtheorem{definition}[theorem]{Definition}
\newtheorem{remark}[theorem]{Remark}
\newtheorem*{remarknn}{Remark}

\newcommand{\R}{\mathbb{R}}
\newcommand{\Z}{\mathbb{Z}}
\newcommand{\C}{\mathbb{C}}
\newcommand{\T}{\mathbb{T}}
\newcommand{\E}{\mathbb{E}}
\newcommand{\PP}{\mathbb{P}}
\newcommand{\eps}{\varepsilon}
\newcommand{\U}{\mathcal U}
\newcommand{\ip}[2]{\langle #1,#2\rangle}
\newcommand{\vvvert}{|\!|\!|}
\DeclareMathOperator{\cth}{coth}
\DeclareMathOperator{\sh}{sinh}
\newcommand{\ch}{\cosh}

\title{Gate (II.11), phase V2:\\[4pt]
\large The interaction comparison on the space--time torus,\\
the renormalization dictionary, and uniform Nelson bounds}
\author{}
\date{16 August 2026}

\begin{document}
\maketitle

\begin{abstract}
This document executes phase V2 of the thermal two-torus route: Statements A4
and A5 of the scoping note are proved. A single Gaussian family realizes the
lattice and continuum periodic free fields on one probability space (the
same-noise, different-multiplier coupling); the two interactions, written in
ensemble-Wick plus finite-recombination form, are compared in $L^2$ with the
explicit rate $O(\eps_N(1+\log(1/\eps_N))^{3/2})$; principal (dispersion +
missing-mode) and aliasing contributions to the squared norm are all
$O(\eps_N^2\log^3(1/\eps_N))$. The renormalization dictionary is proved with
a rate: $\delta_N=c_{N,L}-c^{\mathrm{ct}}_{L,N}
=\tfrac1{2\pi}\log\tfrac4\pi+O(\eps_N)$; together with the equal-point
variance comparison this locks the gate's continuum polynomial to plain
$\varphi^4$, $\delta_\infty$ entering only as the offset between the two
tadpoles and not as a mass counterterm (Corollary~\ref{cor:lock},
Remark~\ref{rem:deltarole}). The mean-shift interpolant of
V1 is identified in closed form,
$H_k(s)=\tfrac{\rho_k}2\,\sh(\gamma_N(k)(s-t/2))/\sh(\gamma_N(k)t/2)$, whence
the shift fields are bounded uniformly in $N$ by coarse data alone and their
symbols converge at the rate $2^{-r(1+\eta)/2}$ from the $D4$ envelope and an
exact filter-tail estimate.  The 16 August repair retains the previously
omitted four-edge sector, separates folded lattice, exact-output mixed, and
continuum shift symbols, takes the real part of the asymmetric mixed surrogate,
and writes the mixed alias term explicitly; no
unnamed remainder is used. Finally, the uniform-in-$N$ Nelson bounds
$\sup_N\|e^{-\U_N}\|_{L^p}<\infty$ (plain and shifted) are proved by the
Part~I Appendix-C pattern, with the comparison scale selected through
$1+\log(1/\eps_{N_b})\approx(b/2c_1t)^{1/2}$ from the depth of the negative
tail and the $L^2$ comparison of this document replacing the mollifier
estimate. Every import is isolated; all constants are explicit in their
dependence on $(m,L,t,\lambda,t_0)$.
\end{abstract}

\tableofcontents

\section{The coupling}\label{sec:coupling}

Fix $L,m>0$, $t\ge t_0>0$ and, after enlarging $t_0$ if necessary,
\begin{equation}
 t_0L\ \ge\ 1 ,
 \tag{$\star$}\label{eq:tLstar}
\end{equation}
which is what Lemma~\ref{lem:envelope}(ii) requires and is assumed throughout this
document; every statement below inherits it. It is not cosmetic --- see the remark
following Theorem~\ref{thm:A4a} --- and it costs nothing, $L$ being fixed and $t\ge t_0$
throughout Part~II. Dual lattices:
\[
 D_N:=\tfrac{2\pi}t\Z\times\Gamma_N,\qquad
 D_\infty:=\tfrac{2\pi}t\Z\times\Gamma_\infty,\qquad
 \Gamma_N=\{\pi j/L:-r_N\le j<r_N\},\quad
 \Gamma_\infty=\tfrac\pi L\Z .
\]
Elements are written $\kappa=(k_0,k)$, $|\kappa|^2=k_0^2+k^2$. Symbols:
\begin{equation}
 \hat c_N(\kappa):=\frac1{k_0^2+\gamma_N(k)^2}\ (\kappa\in D_N),\qquad
 \hat c_\infty(\kappa):=\frac1{k_0^2+\gamma(k)^2},\qquad
 \hat c_\times(\kappa):=\bigl(\hat c_N\hat c_\infty\bigr)^{1/2}(\kappa),
 \label{eq:symbols}
\end{equation}
with $\gamma(k)^2=m^2+k^2$ and $\gamma_N(k)^2=m^2+4\eps_N^{-2}
\sin^2(\eps_Nk/2)$. Throughout, $\log_N:=1+\log(1/\eps_N)$.

\begin{definition}[same-noise coupling]\label{def:coupling}
On a probability space $(\Omega,\PP)$ let $\{G(\kappa)\}_{\kappa\in D_\infty}$
be jointly Gaussian, centered, with
$G(-\kappa)=\overline{G(\kappa)}$ and
$\E[G(\kappa)\overline{G(\kappa')}]=\delta_{\kappa\kappa'}$. Define, for
$z=(s,x)\in\T_t\times\T_L$ and $x\in\Lambda_N$ respectively,
\[
 \eta(z):=(2Lt)^{-1/2}\!\!\sum_{\kappa\in D_\infty}\!\!
 G(\kappa)\,\hat c_\infty(\kappa)^{1/2}e^{i\kappa\cdot z},
 \qquad
 \eta_N(z):=(2Lt)^{-1/2}\!\!\sum_{\kappa\in D_N}\!\!
 G_N(\kappa)\,\hat c_N(\kappa)^{1/2}e^{i\kappa\cdot z},
\]
where $G_N:=G$ off the \emph{edge column} $k_e:=-\pi/\eps_N$, and on it
\[
 G_N(k_0,k_e):=2^{-1/2}\bigl(G(k_0,k_e)+G(k_0,-k_e)\bigr),
 \qquad k_0\in\tfrac{2\pi}t\Z .
\]
The two fields are understood, respectively, as
the first an $L^2(\PP)$-valued distribution (its integrals against
space--time test functions and all Wick integrals below are the objects used;
no pointwise continuum field is needed), the second as an honest jointly
measurable random field (finitely many spatial modes; each mode has a
continuous modification in $s$ by Kolmogorov--Chentsov, since
$\E|X_s-X_{s'}|^2=2(c_\omega(0)-c_\omega(s-s'))\le|s-s'|$ by the Lipschitz
bound $|c_\omega'|\le\tfrac12$, and Gaussianity upgrades this to
$2n$-th moments).
\end{definition}

\begin{remarknn}[why the edge column must be symmetrized]
The band is \emph{asymmetric}: $k_e=-\pi/\eps_N\in\Gamma_N$ while
$-k_e=+\pi/\eps_N\in\Gamma_\infty\setminus\Gamma_N$, so $D_N\ne-D_N$ and the Hermitian
constraint $G(-\kappa)=\overline{G(\kappa)}$ does \emph{not} close on $D_N$ at that one
column. Had $\eta_N$ been defined with $G$ in place of $G_N$ it would not be real:
conjugating its defining sum reindexes it over $-D_N$, every column with
$|k|<\pi/\eps_N$ cancels, and --- using $e^{ik_ex}=e^{-ik_ex}=(-1)^{x/\eps_N}$ on
$\Lambda_N$ --- one is left with
\[
 \overline{\eta_N}-\eta_N=(2Lt)^{-1/2}(-1)^{x/\eps_N}\sum_{k_0}
 \hat c_N(k_0,k_e)^{1/2}e^{ik_0s}\bigl[G(k_0,-k_e)-G(k_0,k_e)\bigr],
\]
a nondegenerate Gaussian: the two families are independent, so by \eqref{eq:ML} its
variance equals $\frac1{2L}\cth(t\gamma_e/2)/\gamma_e=\eps_N/(4L)+O(\eps_N^3)$, where
$\gamma_e:=\gamma_N(k_e)=(m^2+4\eps_N^{-2})^{1/2}$. The induced defect in the covariance
would be the omission of the whole edge column, i.e.\ $\eps_N/(8L)+O(\eps_N^3)$ at
coincident points --- one power of $\eps_N$ \emph{larger} than every term tracked in
\S\ref{sec:A4a}. The symmetrization of $G_N$ is available exactly because $k_e$ and
$-k_e$ carry the same spatial function on $\Lambda_N$ and the same dispersion,
$\gamma_N(k_e)=\gamma_N(-k_e)$. Reality is load-bearing downstream: the pointwise bound
of Lemma~\ref{lem:lower} is an identity between real numbers, and V4's phase estimate
uses $|e^{ia}-e^{ib}|\le|a-b|$ for real $a,b$.
\end{remarknn}

\begin{proposition}[laws and covariances]\label{prop:laws}
\emph{(i)} $\eta_N$ is real, and its law is that of the V1 periodic lattice
ensemble $\mu^{\rm per}_{N,t}$:
\[
 \E[\eta_N(s,x)\eta_N(s',y)]
 =\frac1{2L}\sum_{k\in\Gamma_N}e^{ik(x-y)}\,c_{\gamma_N(k)}(s-s'),
 \qquad
 c_\omega(\tau)=\frac{\cosh(\omega(\tfrac t2-|\tau|_{\T_t}))}
 {2\omega\,\sh(\omega t/2)} .
\]
\emph{(ii)} With $C_\times(z-z'):=(2Lt)^{-1}\sum_{\kappa\in D_N}
\hat c_\times(\kappa)e^{i\kappa(z-z')}$ (and likewise $C_N$, $C_\infty$ with
$\hat c_N$, $\hat c_\infty$), the mixed covariance is
\[
 \E[\eta_N(z)\eta(z')]=C_\times(z-z')+R_N(z,z'),
\]
where $R_N$ is carried by the edge column alone,
\[
 R_N(z,z')=(2Lt)^{-1}\sum_{k_0}\hat c_\times(k_0,k_e)\,e^{ik_0(s-s')}
 \Bigl[\tfrac1{\sqrt2}\bigl(e^{ik_e(x-y)}+e^{ik_e(x+y)}\bigr)-e^{ik_e(x-y)}\Bigr],
\]
and, since the bracket has modulus $\le\sqrt2+1<3$ and, on that column,
$\hat c_N=(k_0^2+\gamma_e^2)^{-1}$ while
$\hat c_\infty=(k_0^2+m^2+\pi^2\eps_N^{-2})^{-1}\le(k_0^2+\gamma_e^2)^{-1}$
(because $\pi^2>4$), whence $\hat c_\times\le(k_0^2+\gamma_e^2)^{-1}$,
\[
 |R_N|\le3(2Lt)^{-1}\!\sum_{k_0}\frac1{k_0^2+\gamma_e^2}
 =\frac3{2L}\,\frac{\cth(t\gamma_e/2)}{2\gamma_e}
 \le\frac{3\,\cth(2t_0)}{8L}\,\eps_N
\]
by \eqref{eq:ML}, using $\gamma_e\ge2\eps_N^{-1}$ and $t/\eps_N\ge2t_0$ for
$\eps_N\le\tfrac12$, $t\ge t_0$.
\emph{(iii)} The edge column is best recorded not through $R_N$ but through the
equivalent \emph{symmetric} description that $R_N$ encodes. Because
$\hat c_\infty(k_0,k_e)=\hat c_\infty(k_0,-k_e)$ and
$e^{ik_ex}=e^{-ik_ex}$ on $\Lambda_N$, the mixed kernel is
\begin{equation*}
 \E[\eta_N(z)\eta(z')]
 =(2Lt)^{-1}\!\!\sum_{\kappa\in\tilde D_N}\!\!
 \chi(k)\,\hat c_\times(\kappa)\,e^{i\kappa(z-z')},
 \qquad
 \tilde D_N:=\tfrac{2\pi}t\Z\times\{\pi j/L:|j|\le r_N\},
\end{equation*}
for $z\in\T_t\times\Lambda_N$, $z'\in\T_t\times\T_L$, with $\chi(k)=1$ for
$|k|<\pi/\eps_N$, $\chi(\pm\pi/\eps_N)=2^{-1/2}$, and
$\hat c_\times(k_0,+\pi/\eps_N):=\hat c_\times(k_0,k_e)$: the mixed pairing runs over
the \emph{symmetric} band $\tilde D_N$, each edge momentum carrying weight $2^{-1/2}$.
(In terms of the noise families, the edge column has \emph{two} partners, not one:
$\E[G_N(k_0,k_e)G(\kappa')]
=2^{-1/2}\bigl(\delta_{\kappa',(-k_0,\pi/\eps_N)}+\delta_{\kappa',(-k_0,-\pi/\eps_N)}\bigr)$,
which is what produces both the $e^{ik_e(x-y)}$ and the $e^{ik_e(x+y)}$ term of $R_N$.)
Consequently:
\emph{(a)} in the second-moment sums the two edge configurations
$(\pm\pi/\eps_N,\mp\pi/\eps_N)$ each carry $\chi^2=\tfrac12$ and together restore
\emph{full} weight, so those sums are \emph{unchanged};
\emph{(b)} In the fourth-moment cross sum, let
\begin{equation}
 \mathcal E_N^{\rm edge}:=4!(2Lt)^{-2}
 \left[
 \sum_{\substack{\kappa_i\in\tilde D_N^4\\\sum\kappa_i=0}}
       \prod_i\chi(k_i)\hat c_\times(\kappa_i)
 -\sum_{\substack{\kappa_i\in D_N^4\\\sum\kappa_i=0}}
       \prod_i\hat c_\times(\kappa_i)
 \right].
 \label{eq:Eedge}
\end{equation}
Partition the difference by the number $|J|$ of edge momenta.  Exactly three edge
momenta are impossible: their sum is $\pm\pi/\eps_N$ or
$\pm3\pi/\eps_N$, so cancellation by a non-edge fourth momentum would require an
edge momentum or a momentum outside the band.  Four edge momenta are possible and
must be retained.

\emph{One edge factor:} it contributes
$\sum_{k_0}\hat c_\times(k_0,k_e)\le\tfrac{t\eps_N}4\cth(2t_0)$, while the
complementary three-fold sum is at most $F_3$ at spatial momentum $\pi/\eps_N$, i.e.\
$\le C(m,t_0)(tL)^2\eps_N^2\log_N^2$ by \eqref{eq:convdecay}; after the
normalizing prefactor the total is $O(t\eps_N^3\log_N^2)$.

\emph{Two edge factors:} the discrepant type is
$(+\pi/\eps_N,-\pi/\eps_N)$. With
$f(k_0):=\hat c_\times(k_0,k_e)$ this class is
$\le(2Lt)^{-2}\binom42\|f*f\|_\infty\sum_wF_2(0,-w)$, and
$\|f*f\|_\infty\le\sum_{k_0}f(k_0)^2\le\sum_{k_0}(k_0^2+\gamma_e^2)^{-2}
=O(t\eps_N^3)$, giving $O(t\eps_N^3)$ after normalization.

\emph{Four edge factors:} exact spatial conservation requires two momenta
$+\pi/\eps_N$ and two momenta $-\pi/\eps_N$.  There are $\binom42=6$
orders and each has weight $(2^{-1/2})^4=1/4$.  Their contribution is
\begin{equation}
 \Delta_{4e}:=4!(2Lt)^{-2}\frac32
 \sum_{k_{0,1}+\cdots+k_{0,4}=0}\prod_{i=1}^4f(k_{0,i}).
 \label{eq:four-edge}
\end{equation}
Put $A:=\sum_{k_0}f(k_0)$ and $B:=\sum_{k_0}f(k_0)^2$.  Since
$f(k_0)\le(k_0^2+\gamma_e^2)^{-1}$, $\gamma_e\ge2\eps_N^{-1}$, and
$t/\eps_N\ge2t_0$,
\[
 A\le\frac{t\eps_N}{4}\cth(2t_0),\qquad
 B\le\gamma_e^{-2}A\le\frac{t\eps_N^3}{16}\cth(2t_0).
\]
For convolution on $(2\pi/t)\Z$,
\[
 \sum_{\sum_i k_{0,i}=0}\prod_i f(k_{0,i})
 =\|f*f\|_2^2
 \le\|f*f\|_\infty\|f*f\|_1
 \le BA^2.
\]
Consequently
\begin{equation}
 0\le\Delta_{4e}\le
 \frac{9t}{256L^2}\cth(2t_0)^3\eps_N^5.
 \label{eq:four-edge-bound}
\end{equation}
Combining the three possible sectors gives
\begin{equation}
 |\mathcal E_N^{\rm edge}|
 \le C(m,t_0,L)t\eps_N^3\log_N^2+\Delta_{4e}.
 \label{eq:Eedge-bound}
\end{equation}
This replaces the false assertion that all sectors with $|J|\ge3$ are impossible.
\end{proposition}

\begin{proof}
\emph{Reality of $\eta_N$.} Off the edge column both $\kappa$ and $-\kappa$ lie in $D_N$
and $G_N(-\kappa)=\overline{G_N(\kappa)}$, $\hat c_N(-\kappa)=\hat c_N(\kappa)$, so those
terms pair into a real sum. On the edge column the relevant involution of $D_N$ is
$(k_0,k_e)\mapsto(-k_0,k_e)$, because $e^{ik_ex}=e^{-ik_ex}$ for $x\in\Lambda_N$; and
by the edge-column convention of Definition~\ref{def:coupling},
\[
 \begin{split}
 G_N(-k_0,k_e)&=2^{-1/2}\bigl(G(-k_0,k_e)+G(-k_0,-k_e)\bigr)\\
 &=2^{-1/2}\overline{\bigl(G(k_0,-k_e)+G(k_0,k_e)\bigr)}
 =\overline{G_N(k_0,k_e)},
 \end{split}
\]
while $\hat c_N(-k_0,k_e)=\hat c_N(k_0,k_e)$. Hence that column is real too, and
$\eta_N$ is real. (Without the symmetrization of $G_N$ this step fails; see the remark following
Definition~\ref{def:coupling}.)

\emph{Orthonormality of $G_N$ on $D_N$.} $\E[G_N(\kappa)\overline{G_N(\kappa')}]
=\delta_{\kappa\kappa'}$: off the edge this is the hypothesis, on the edge
$\tfrac12(\delta+0+0+\delta)=\delta_{k_0k_0'}$, and an edge/non-edge pair is orthogonal
since $\pm k_e\ne k$ for $k\in\Gamma_N\setminus\{k_e\}$. Likewise, using
$\E[G(\kappa)G(\kappa')]=\delta_{\kappa,-\kappa'}$,
\[
 \E\bigl[G_N(k_0,k_e)G_N(k_0',k_e)\bigr]
 =\tfrac12\bigl(0+\delta_{k_0,-k_0'}+\delta_{k_0,-k_0'}+0\bigr)=\delta_{k_0,-k_0'},
\]
so the edge column contributes to $\E[\eta_N(z)\eta_N(z')]$ the term
\[
 (2Lt)^{-1}\sum_{k_0}\hat c_N(k_0,k_e)\,e^{ik_0(s-s')}e^{ik_e(x+y)},
\]
and $e^{ik_e(x+y)}=e^{ik_e(x-y)}$ on $\Lambda_N$: exactly the $k=k_e$ term of
the displayed kernel, with full weight. Off the edge the pairing
$\kappa'=-\kappa$ is the usual one.
The covariance is therefore as stated; for (i) it remains to evaluate the $k_0$-sum:
\begin{equation}
 \frac1t\sum_{k_0\in\frac{2\pi}t\Z}\frac{e^{ik_0\tau}}{k_0^2+\omega^2}
 =c_\omega(\tau)\qquad(\omega>0),
 \label{eq:ML}
\end{equation}
which holds because both sides are continuous in $\tau$, the right side has
Fourier coefficients $t^{-1}\int_0^te^{-ik_0\tau}c_\omega(\tau)d\tau
=t^{-1}(k_0^2+\omega^2)^{-1}$ --- integrate
$(-\partial_\tau^2+\omega^2)c_\omega=\delta_{\T_t}$ against $e^{-ik_0\tau}$,
where the distributional identity on the circle follows from
$c_\omega''=\omega^2c_\omega$ off $0$ and the derivative jump
$c_\omega'(0^+)-c_\omega'(0^-)=-1$ (direct differentiation) --- and the left
side is the absolutely convergent Fourier series of a continuous function.
This matches V1, Remark 5.5, so the laws agree.
\end{proof}

\section{Wick structure and the two interactions}\label{sec:wick}

Equal-point variances (finite on the lattice):
\[
 v_N(t):=\E[\eta_N(z)^2]
 =\frac1{2L}\sum_{k\in\Gamma_N}\frac{\cth(t\gamma_N(k)/2)}{2\gamma_N(k)},
 \qquad
 c_{N,L}=\frac1{2L}\sum_{k\in\Gamma_N}\frac1{2\gamma_N(k)} .
\]

\begin{lemma}[recombination; the finite constants $d$]\label{lem:recomb}
For any real random variable $X$ with variance $v$ and any constant $c$,
\begin{equation}
 {:}X^4{:}_c={:}X^4{:}_v+6(v-c)\,{:}X^2{:}_v+3(v-c)^2,
 \qquad
 {:}X^2{:}_c={:}X^2{:}_v+(v-c),
 \label{eq:recomb}
\end{equation}
where ${:}X^j{:}_a$ denotes Wick ordering with constant $a$
(${:}X^2{:}_a=X^2-a$, ${:}X^4{:}_a=X^4-6aX^2+3a^2$). Moreover
\begin{equation}
 d_N(t):=v_N(t)-c_{N,L}
 =\frac1{2L}\sum_{k\in\Gamma_N}
 \frac{\cth(t\gamma_N(k)/2)-1}{2\gamma_N(k)}
 =\frac1{2L}\sum_{k\in\Gamma_N}
 \frac1{\gamma_N(k)\,(e^{t\gamma_N(k)}-1)},
 \label{eq:dN}
\end{equation}
and, with $d_\infty(t)$ the same sum over $\Gamma_\infty$ with $\gamma$ and
$\gamma_-(k):=\max(m,\tfrac2\pi|k|)$:
\[
 0\ \le\ d_N(t)\ \le\ \bar d:=\frac1{2L}\sum_{k\in\Gamma_\infty}
 \frac1{\gamma_-(k)\bigl(e^{t_0\gamma_-(k)}-1\bigr)}\ <\ \infty
 \qquad\text{uniformly in $N$ and $t\ge t_0$},
\]
and $|d_N(t)-d_\infty(t)|\le C(m,L,t_0)\,\eps_N^2$, uniformly in $t\ge t_0$.
\end{lemma}

\begin{proof}
\eqref{eq:recomb} is polynomial algebra:
$X^4-6cX^2+3c^2=(X^4-6vX^2+3v^2)+6(v-c)(X^2-v)+3(v-c)^2$. In \eqref{eq:dN},
$\cth(y/2)-1=2/(e^y-1)$. Uniform bound: on the band,
$\gamma_N(k)\ge\gamma_-(k)$ (from $\sin u\ge\frac2\pi u$ on
$[0,\pi/2]$), $y\mapsto(y(e^{ty}-1))^{-1}$ is decreasing in both $y$ and $t$,
and $\Gamma_N\subset\Gamma_\infty$ with positive summands, so $d_N(t)\le\bar
d$; $\bar d<\infty$ since its summands decay like $e^{-2t_0|k|/\pi}$.
Convergence with rate: for fixed $k$, $f(y):=(y(e^{ty}-1))^{-1}$ has
\begin{gather*}
 |f'(y)|=\frac{(e^{ty}-1)+ty\,e^{ty}}{y^2(e^{ty}-1)^2}
 \le\frac1{y^2(e^{ty}-1)}+\frac{t\,e^{ty}}{y(e^{ty}-1)^2}
 \le C\,(1+t)\,\frac{e^{-ty}}{\min(y,1)^2}\\
 (y\ge m,\ ty\ge t_0m),
\end{gather*}
and $0\le\gamma(k)-\gamma_N(k)\le\eps_N^2k^4/(24m)$ (free-rate document,
Lemma~3.1, eq.~(3.1)); the mean value theorem on
$[\gamma_N(k),\gamma(k)]\subset[\gamma_-(k),\infty)$ gives per-$k$ errors
$\le C(m)\,(1+t)e^{-t\gamma_-(k)}\eps_N^2k^4$, whose $\frac1{2L}$-weighted sum
is $\le C(m,L,t_0)\eps_N^2$ uniformly in $t\ge t_0$ (as
$\sup_{t\ge t_0}(1+t)e^{-t\gamma_-/2}<\infty$ and
$\sum_ke^{-t_0\gamma_-(k)/2}k^4<\infty$); the missing modes
$k\in\Gamma_\infty\setminus\Gamma_N=\{k\ge\pi/\eps_N\}\cup\{k<-\pi/\eps_N\}$ --- note
that $+\pi/\eps_N$ \emph{is} missing while $-\pi/\eps_N$ is not, the band being
asymmetric --- contribute $\le Ce^{-2t_0/\eps_N}$.
\end{proof}

\begin{definition}[the interactions, ensemble-Wick form]\label{def:U}
With $\mathrm{Vol}:=2Lt$:
\begin{align}
 \U_N&:=\lambda\Bigl[X_4+6\,d_N(t)\,X_2+3\,d_N(t)^2\,\mathrm{Vol}\Bigr],&
 X_j&:=\eps_N\!\!\sum_{x\in\Lambda_N}\int_0^t
 {:}\eta_N(s,x)^j{:}_{v_N}\,ds,
 \label{eq:UN}\\
 \U&:=\lambda\Bigl[Y_4+6\,d_\infty(t)\,Y_2+3\,d_\infty(t)^2\,
 \mathrm{Vol}\Bigr],&
 Y_j&:=\int_0^t\!\!\int_{\T_L}{:}\eta(s,x)^j{:}\,dx\,ds ,
 \label{eq:Uc}
\end{align}
where $Y_j$ is defined as the $L^2(\PP)$ limit of its spatial mode
truncations $Y_j^{(K)}$ (Proposition~\ref{prop:Yexists}). By
Lemma~\ref{lem:recomb}, $\U_N$ equals the V1 interaction
$\lambda\eps_N\sum_x\int_0^t{:}\eta_N^4{:}_{c_{N,L}}ds$ exactly; $\U$ is the
continuum interaction completely Wick-ordered by the continuum field's own
vacuum constant --- the truncation convention; by Corollary~\ref{cor:lock}
this coincides in the limit with the lattice-vacuum convention, with no
counterterm between them.
\end{definition}

\begin{lemma}[Wick--Isserlis pairings]\label{lem:isserlis}
If $(X,Y)$ are jointly Gaussian, centered, with variances $v_X,v_Y$ and
covariance $c=\E[XY]$, then for $j,l\ge0$,
\[
 \E\bigl[{:}X^j{:}_{v_X}\,{:}Y^l{:}_{v_Y}\bigr]=\delta_{jl}\;j!\;c^{\,j}.
\]
\end{lemma}

\begin{proof}
From $\E[e^{\alpha X+\beta Y}]
=e^{(\alpha^2v_X+2\alpha\beta c+\beta^2v_Y)/2}$ and
${:}e^{\alpha X}{:}:=e^{\alpha X-\alpha^2v_X/2}$ we get
\[
 \E\bigl[{:}e^{\alpha X}{:}\,{:}e^{\beta Y}{:}\bigr]=e^{\alpha\beta c};
\]
expand both
sides (${:}e^{\alpha X}{:}=\sum_j\alpha^j{:}X^j{:}/j!$, absolutely convergent
in $L^2$) and compare coefficients of $\alpha^j\beta^l$.
\end{proof}

\begin{proposition}[existence of $Y_j$; positivity of the mode sums]
\label{prop:Yexists}
Let $Y_j^{(K)}$ be defined with $\eta$ replaced by its truncation to spatial
modes $|k|\le K$ (a smooth Gaussian field), Wick-ordered by its own variance.
Then, for $j\in\{2,4\}$,
\[
 \E\bigl[Y_j^{(K)}Y_j^{(K')}\bigr]
 =j!\;(2Lt)^{-(j-2)}
 \sum_{\substack{\kappa_1,\dots,\kappa_j\in D_\infty^{\le K\wedge K'}\\
 \kappa_1+\dots+\kappa_j=0}}\;\prod_{i=1}^j\hat c_\infty(\kappa_i)
 =:j!\,(2Lt)^{-(j-2)}\,S_j^{(K\wedge K')},
\]
$S_j^{(K)}$ is nondecreasing in $K$ and bounded
(Lemma~\ref{lem:envelope}\,(iii)), hence $Y_j^{(K)}$ is $L^2(\PP)$-Cauchy and
$Y_j:=\lim Y_j^{(K)}$ exists, with
$\E[Y_j^2]=j!(2Lt)^{-(j-2)}S_j^{(\infty)}$.
\end{proposition}

\begin{proof}
By Lemma~\ref{lem:isserlis},
$\E[Y_j^{(K)}Y_j^{(K')}]=j!\iint C_{K\wedge K'}(z-z')^j\,dz\,dz'$, where
$C_{K\wedge K'}$ is the covariance of the two truncations (Wick pairings only
connect modes present in both fields, and the joint covariance of the two
truncated fields is the min-truncated kernel). Expanding $C^j$ in its
absolutely convergent Fourier series ($C^j$ carries the prefactor
$(2Lt)^{-j}$) and integrating over $(\T_t\times\T_L)^2$ gives the displayed
sum ($\iint_{(\T_t\times\T_L)^2}
e^{i\Sigma\kappa\cdot z}e^{-i\Sigma\kappa\cdot z'}\,dz\,dz'
=(2Lt)^2$ if $\Sigma\kappa=0$, else $0$, so the net power is
$(2Lt)^{-(j-2)}$; absolute convergence of the interchange is
Lemma~\ref{lem:envelope}\,(iii)). All summands are positive, so
$K\mapsto S_j^{(K)}$ is nondecreasing and bounded, hence convergent; then
$\|Y^{(K)}-Y^{(K')}\|_2^2
=j!(2Lt)^{-(j-2)}(S^{(K)}+S^{(K')}-2S^{(K\wedge K')})\to0$.
\end{proof}

\section{Covariance estimates}\label{sec:covest}

\begin{lemma}[envelope, decay of convolutions, summability]\label{lem:envelope}
Set $\hat c_{\mathrm{env}}(\kappa):=\dfrac{\pi^2/4}{|\kappa|^2+m^2}$ on
$D_\infty$. Then:
\emph{(i)} $\hat c_N\le\hat c_{\mathrm{env}}$ on $D_N$,
$\hat c_\infty\le\hat c_{\mathrm{env}}$, and
$\hat c_\times\le\hat c_{\mathrm{env}}$ on $D_N$.
\emph{(ii)} With $\Sigma_1:=\sum_{\kappa\in D_\infty}
\hat c_{\mathrm{env}}(\kappa)^2<\infty$ and the two- and three-fold
convolutions (sums over $D_\infty$)
\[
 F_2(\chi):=\sum_{\kappa}\hat c_{\mathrm{env}}(\kappa)
 \hat c_{\mathrm{env}}(\chi-\kappa),
 \qquad
 F_3:=\hat c_{\mathrm{env}}*F_2 :
\]
\begin{equation}
 F_2(\chi)\le C_2\,tL\;\frac{1+\log(2+|\chi|)}{m^2+|\chi|^2},
 \qquad
 F_3(\chi)\le C_3\,(tL)^2\,\frac{(1+\log(2+|\chi|))^2}{m^2+|\chi|^2},
 \label{eq:convdecay}
\end{equation}
with $C_2,C_3$ depending on $(m,t_0)$ (the isolated temporal-grid term uses
$t\ge t_0$); $tL\ge1$ is also assumed so that the isolated two-dimensional
grid term is absorbed by the density term.  No later statement treats these
constants as functions of $m$ alone.
\emph{(iii)} $S_4^{(\infty)}\le\sum_{\kappa}\hat c_{\mathrm{env}}(\kappa)
F_3(\kappa)<\infty$ and
$S_2^{(\infty)}\le\Sigma_1<\infty$; the same bounds hold for all the mixed
and lattice sums below, aliased or not, since every symbol involved is
dominated by $\hat c_{\mathrm{env}}$ at its own argument.
\end{lemma}

\begin{proof}
(i) On the band, $\sin u\ge\tfrac2\pi u$ for $u\in[0,\pi/2]$ gives
$4\eps^{-2}\sin^2(\eps k/2)\ge(4/\pi^2)k^2$, so
$k_0^2+\gamma_N^2\ge k_0^2+m^2+(4/\pi^2)k^2
\ge(4/\pi^2)(|\kappa|^2+m^2)$; the $\hat c_\times$ bound follows since a
geometric mean of two numbers $\le a$ is $\le a$.
(ii) $\Sigma_1<\infty$: the grid $D_\infty$ has cell area
$(2\pi/t)(\pi/L)$ and $\int_{\R^2}(|\kappa|^2+m^2)^{-2}d\kappa<\infty$; the
standard sum--integral comparison (the summand is monotone along rays outside
a bounded set) gives $\Sigma_1\le C(m)\,tL$. Throughout we use the two grid
counts (same comparison; density $\tfrac{tL}{2\pi^2}$ per unit area):
\begin{equation}
 \begin{split}
 \sum_{|\kappa|\le R}\hat c_{\mathrm{env}}(\kappa)
 &\le c_0(m)\,tL\,(1+\log(2+R)),\\
 \sum_{|\kappa|>R}\frac{1+\log(2+|\kappa|)}{(m^2+|\kappa|^2)^2}
 &\le c_0(m)\,tL\,\frac{1+\log(2+R)}{m^2+R^2},
 \end{split}
 \label{eq:gridcounts}
\end{equation}
and the monotone majorization: for $R\ge|\chi|/2$,
$\frac{1+\log(2+R)}{m^2+R^2}\le c_1(m)\frac{1+\log(2+|\chi|)}{m^2+|\chi|^2}$
(the function is decreasing beyond a bounded set, and on the bounded set both
sides are comparable to constants).
\emph{$F_2$:} split $D_\infty$ into
$\mathrm I=\{|\kappa|\le|\chi|/2\}$,
$\mathrm{II}=\{|\chi-\kappa|\le|\chi|/2\}$, and the remainder $\mathrm{III}$
(on which both $|\kappa|,|\chi-\kappa|>|\chi|/2$). On $\mathrm I$:
$|\chi-\kappa|\ge|\chi|/2$, so
$\hat c_{\mathrm{env}}(\chi-\kappa)\le c_1\frac{\pi^2/4}{m^2+|\chi|^2}$
and the first count in \eqref{eq:gridcounts} applies to
$\sum_{\mathrm I}\hat c_{\mathrm{env}}(\kappa)$. On $\mathrm{II}$:
symmetric. On $\mathrm{III}$: for each $\kappa$ let
$w(\kappa)$ be whichever of $\kappa,\chi-\kappa$ has the \emph{larger}
modulus and $w'(\kappa)$ the other; then $|w|\ge|w'|>|\chi|/2$, and since
$R\mapsto\frac{\pi^2/4}{m^2+R^2}$ is decreasing, the product of the two
factors is at most $\bigl(\frac{\pi^2/4}{m^2+|w'|^2}\bigr)^2$. At most two
$\kappa$ produce the same $w'$, so by the second count in
\eqref{eq:gridcounts} (without the logarithm),
$\sum_{\mathrm{III}}\le2\sum_{|w'|>|\chi|/2}
\bigl(\tfrac{\pi^2/4}{m^2+|w'|^2}\bigr)^2
\le C(m)\,tL\,(m^2+|\chi|^2)^{-1}$. Summing the three regions
gives the $F_2$ bound.
\emph{$F_3=\hat c_{\mathrm{env}}*F_2$:} same three regions with
$f=F_2$, $g=\hat c_{\mathrm{env}}$. Region $\mathrm I$
($|u|\le|\chi|/2$, $u$ the $F_2$-argument):
$g(\chi-u)\le c_1\frac{\pi^2/4}{m^2+|\chi|^2}$ and
$\sum_{|u|\le|\chi|/2}F_2(u)\le C_2tL\sum_{|u|\le|\chi|/2}
\frac{1+\log(2+|u|)}{m^2+|u|^2}\le C(m)(tL)^2(1+\log(2+|\chi|))^2$
(one logarithm from the summand, one from the radial integral). Region
$\mathrm{II}$ ($|\chi-u|\le|\chi|/2$, so $|u|\ge|\chi|/2$):
$F_2(u)\le C(m)tL\,c_1\frac{1+\log(2+|\chi|)}{m^2+|\chi|^2}$ by monotone
majorization, and $\sum_{|\chi-u|\le|\chi|/2}\hat c_{\mathrm{env}}(\chi-u)
\le c_0tL(1+\log(2+|\chi|))$. Region $\mathrm{III}$ (both moduli
$>|\chi|/2$): with $w$ the larger-modulus argument as before, majorize the
$w$-factor's decay profile at $|w'|$ (monotone majorization, valid since
$|w|\ge|w'|$) so that the product is dominated by
$C(m)tL\,\frac{1+\log(2+|w'|)}{(m^2+|w'|^2)^2}$, and apply the second count
in \eqref{eq:gridcounts} over $\{|w'|>|\chi|/2\}$:
$\mathrm{III}\le C(m)(tL)^2\frac{1+\log(2+|\chi|)}{m^2+|\chi|^2}$. Summing
the regions gives the $F_3$ bound.
(iii) $S_4^{(\infty)}\le\sum_\kappa\hat c_{\mathrm{env}}(\kappa)F_3(-\kappa)
\le C(tL)^2\sum_\kappa\frac{(1+\log(2+|\kappa|))^2}{(m^2+|\kappa|^2)^2}
<\infty$.
\end{proof}

\begin{lemma}[symbol differences]\label{lem:symdiff}
For $\kappa\in D_N$:
\begin{equation}
 \bigl|\hat c_N(\kappa)-\hat c_\infty(\kappa)\bigr|
 \le\frac{\eps_N^2k^4}{12}\,\hat c_N(\kappa)\hat c_\infty(\kappa)
 \le\frac{\pi^4}{192}\,\eps_N^2k^4\,
 \hat c_{\mathrm{env}}(\kappa)\,
 \frac{1}{m^2+|\kappa|^2},
 \label{eq:symdiff}
\end{equation}
and $|\hat c_\times-\hat c_\infty|\le|\hat c_N-\hat c_\infty|$ (from
$|\sqrt a-\sqrt b|=\frac{|a-b|}{\sqrt a+\sqrt b}$ and
$\sqrt{\hat c_\infty}\le\sqrt{\hat c_N}+\sqrt{\hat c_\infty}$).
\end{lemma}

\begin{proof}
$\hat c_N-\hat c_\infty
=(\gamma^2-\gamma_N^2)\hat c_N\hat c_\infty$ and
$0\le k^2-4\eps^{-2}\sin^2(\eps k/2)\le\eps^2k^4/12$ (free-rate document,
(3.5)); then insert the envelope twice.
\end{proof}

\section{The fourth-moment comparison (A4(a))}\label{sec:A4a}

\begin{lemma}[the three sums]\label{lem:threesums}
With $X_4,Y_4$ as in Definition~\ref{def:U},
\begin{align}
 \E[X_4^2]&=4!\,(2Lt)^{-2}
 \sum_{\substack{\kappa_i\in D_N^4,\ \Sigma k_0=0\\
 \Sigma k\equiv0\ (\mathrm{mod}\ 2\pi/\eps_N)}}
 \prod_{i=1}^4\hat c_N(\kappa_i),
 \label{eq:XX}\\
 \E[Y_4^2]&=4!\,(2Lt)^{-2}
 \sum_{\kappa_i\in D_\infty^4,\ \Sigma\kappa=0}
 \prod_{i=1}^4\hat c_\infty(\kappa_i).
 \nonumber
\end{align}
Define the asymmetric-band surrogate
\[
 S_{\times,N}^{D}:=4!\,(2Lt)^{-2}
 \sum_{\kappa_i\in D_N^4,\ \Sigma\kappa=0}
 \prod_{i=1}^4\hat c_\times(\kappa_i).
\]
Proposition~\ref{prop:laws}(iii) says that the actual mixed pairing runs over
the \emph{symmetric} band
$\tilde D_N=\tfrac{2\pi}t\Z\times\{\pi j/L:|j|\le r_N\}$ with edge weight
$\chi$.  Therefore the exact cross identity is
\[
 \E[X_4Y_4]
 =4!\,(2Lt)^{-2}\!\!\sum_{\kappa_i\in\tilde D_N^4,\ \Sigma\kappa=0}
 \prod_{i=1}^4\chi(k_i)\hat c_\times(\kappa_i)
 =S_{\times,N}^{D}+\mathcal E_N^{\rm edge}.
\]
Thus the asymmetric sum is never asserted to be the true cross moment.  The bound
\eqref{eq:Eedge-bound}, including the four-edge term
\eqref{eq:four-edge}, is carried explicitly into the proof of
Theorem~\ref{thm:A4a}. In the second-chaos analogue
the two descriptions agree \emph{exactly}, by Proposition~\ref{prop:laws}(iii)(a).
The lattice constraint splits into the \emph{principal} part $\Sigma k=0$ and
the \emph{alias} part $\Sigma k\in\{\pm2\pi/\eps_N,-4\pi/\eps_N\}$
(no other values occur, since each $k_i\in[-\pi/\eps_N,\pi/\eps_N)$).
\end{lemma}

\begin{proof}
By Lemma~\ref{lem:isserlis}, $\E[X_4^2]=4!\,\eps_N^2\sum_{x,x'}
\iint C_N(s-s',x-x')^4$, and
$\eps_N\sum_{u\in\Lambda_N}e^{iKu}=2L\cdot\mathbf 1[K\in(2\pi/\eps_N)\Z]$,
$\int_0^te^{ik_0\tau}d\tau=t\,\mathbf 1[k_0=0]$ applied to the absolutely
convergent (Lemma~\ref{lem:envelope}) Fourier expansion of $C_N^4$; the
$\mathrm{Vol}$-factors combine to $(2Lt)^{2}(2Lt)^{-4}$. For the cross term,
$\int_{\T_L}dy$ enforces exact conservation of the $D_\infty$-labels carried by
$\eta$; by Definition~\ref{def:coupling} those labels run over $\tilde D_N$, the two
edge momenta $\pm\pi/\eps_N$ each carrying the weight $2^{-1/2}\hat c_\times$
(Proposition~\ref{prop:laws}(iii)), and the lattice constraint $\Sigma k\equiv0$ is
then automatic. It is \emph{not} true that modes of $\eta$ outside $D_N$ have zero
covariance with $\eta_N$: the mode at $+\pi/\eps_N$ does not, and that is exactly what
the symmetrization of $G_N$ creates. The alias values:
$\Sigma k\in[-4\pi/\eps_N,4\pi/\eps_N)$ intersected with
$(2\pi/\eps_N)\Z$.
\end{proof}

\begin{lemma}[alias bound]\label{lem:alias}
\begin{equation}
 \mathrm{Al}_N:=4!\,(2Lt)^{-2}\!\!\!
 \sum_{\substack{\kappa_i\in D_N^4,\ \Sigma k_0=0\\
 \Sigma k\in\{\pm2\pi/\eps_N,\,-4\pi/\eps_N\}}}\!\!\!
 \prod_i\hat c_N(\kappa_i)
 \;\le\;C_{\mathrm{al}}(m,t_0)\;tL\;
 \eps_N^{2}\,\bigl(1+\log(1/\eps_N)\bigr)^{3}.
 \label{eq:alias}
\end{equation}
\end{lemma}

\begin{proof}
Parametrize by $(\kappa_1,\kappa_2,\kappa_3)$; then
$\kappa_4=\sigma-\kappa_1-\kappa_2-\kappa_3$ with
$\sigma=(0,\Sigma k)$, $|\sigma|\ge2\pi/\eps_N$. Domination by the envelope
and Lemma~\ref{lem:envelope}(ii):
\[
 \sum_{\kappa_2,\kappa_3}\hat c_{\mathrm{env}}(\kappa_2)
 \hat c_{\mathrm{env}}(\kappa_3)
 \hat c_{\mathrm{env}}(\sigma-\kappa_1-\kappa_2-\kappa_3)
 =F_3(\sigma-\kappa_1)
 \le C_3\,(tL)^2\,\frac{(1+\log(2+|\sigma-\kappa_1|))^2}
 {m^2+|\sigma-\kappa_1|^2}\,,
\]
The \emph{spatial} component of $\sigma-\kappa_1$ has modulus
$\ge2\pi/\eps_N-\pi/\eps_N=\pi/\eps_N$ (and $\ge3\pi/\eps_N$ in the corner
case), and, since $x\mapsto(1+\log(2+x))^2/(m^2+x^2)$ is decreasing for
$x\ge x_0(m)$, the supremum of \eqref{eq:convdecay}'s right-hand side over
$|\sigma-\kappa_1|\ge\pi/\eps_N$ is attained at the endpoint, giving
$F_3(\sigma-\kappa_1)\le C_3(tL)^2\eps_N^2\pi^{-2}
(1+\log(2+4\pi/\eps_N))^2$ uniformly in $\kappa_1$ (for
$\eps_N\le\eps_0(m)$; the finitely many larger $\eps_N$ are absorbed into
$C_{\mathrm{al}}$). Next,
\[
 \sum_{\kappa_1\in D_N}\hat c_{\mathrm{env}}(\kappa_1)
 \le\sum_{k\in\Gamma_N}\sum_{k_0\in\frac{2\pi}t\Z}
 \frac{\pi^2/4}{m^2+k_0^2+k^2}
 \le C\,t\!\!\sum_{k\in\Gamma_N}\!\!\frac1{(m^2+k^2)^{1/2}}
 \le C(m,t_0)\,tL\,\log_N
\]
(one-dimensional sum--integral comparisons in $k_0$, then in $k$ up to the
band radius $\pi/\eps_N$). Multiply the three bounds and the prefactor
$4!(2Lt)^{-2}$; there are three alias values of $\sigma$, absorbed into
$C_{\mathrm{al}}$.
\end{proof}

\begin{theorem}[A4(a): interaction $L^2$ comparison with rate]
\label{thm:A4a}
Assume the standing hypothesis $tL\ge1$ of Lemma~\ref{lem:envelope}(ii). There is
$C=C(m,\lambda,t_0)<\infty$ such that for all $N$ with
$\eps_N\le\tfrac12$ and all $t\ge t_0$:
\begin{equation}
 \bigl\|\,\U_N-\U\,\bigr\|_{L^2(\PP)}
 \;\le\;C\,(1+t)(1+L)\;\eps_N\,\bigl(1+\log(1/\eps_N)\bigr)^{3/2}.
 \label{eq:A4a}
\end{equation}
\end{theorem}

\begin{remarknn}[the r\^ole of $tL\ge1$, and where it is consumed]
The hypothesis is not cosmetic. Lemma~\ref{lem:envelope}(ii) bounds
$\Sigma_1=\sum_{\kappa\in D_\infty}\hat c_{\mathrm{env}}(\kappa)^2$ by $C(m)\,tL$
through a sum--integral comparison on a grid of density $tL/2\pi^2$ plus an isolated
maximal term, and that maximal term is absorbable only when $tL\gtrsim1$. As
$t\downarrow0$ at fixed $L$ the $k_0$-spacing $2\pi/t$ diverges, only the row $k_0=0$
survives, and $\Sigma_1$ \emph{saturates} at a $t$-independent value while $tL\to0$:
numerically ($m=L=1$) $\Sigma_1=7.658,\,6.230,\,6.202,\,6.201,\,6.201$ at
$t=4,1,\tfrac14,\tfrac1{20},\tfrac1{100}$, so $\Sigma_1/(tL)=1.9,\,6.2,\,24.8,\,124,\,620$
grows like $1/t$. Hence \eqref{eq:A4a} must not be invoked at arbitrarily small $t$.
The one place in the chain that would otherwise do so is V4, Theorem 5.2(a); it is
stated there for $t\ge t_*$ with $t_*>0$ fixed, which is all the scoping note's
Laplace-transfer lemma requires (see the amended hypothesis (i) of that lemma).
Throughout Part~II $L$ is fixed and $t\ge t_0$, so $tL\ge1$ costs nothing after
enlarging $t_0$.
\end{remarknn}

\begin{proof}
Throughout, two one-dimensional sum--integral comparisons are used (both
standard: even unimodal summands, density $\tfrac t{2\pi}$ resp.\
$\tfrac L\pi$ plus the maximal term; $j\in\{0,1,2\}$ and any $s>1/2$
(used below at $s=2,3,4$),
$a^2:=m^2+k^2$):
\begin{equation}
 \sum_{k_0\in\frac{2\pi}t\Z}
 \frac{(1+\log(2+|\kappa|))^j}{(m^2+k_0^2+k^2)^{s}}
 \le C_s(m,t_0)\,t\,\frac{(1+\log(2+|k|))^j}{a^{2s-1}},
 \label{eq:k0count}
\end{equation}
(split the sum at $|k_0|\le a$, where $\log(2+|\kappa|)\le
C(1+\log(2+|k|))$, and use
$\int_a^\infty v^{-2s}(1+\log(2+v))^j\,dv\le
C_{j,s}\,a^{1-2s}(1+\log(2+a))^j$ on the rest --- valid for all $a>0$ and
$s>1/2$, splitting at $\max(a,1)$; the isolated maximal term of the near part
is absorbed via $1/a\le1/m\le t/(mt_0)$, which is the sole source of the
$t_0$-dependence of $C_s$), and
\begin{equation}
 \begin{split}
 \sum_{k\in\Gamma_N}\frac{k^4(1+\log(2+|k|))^2}{(m^2+k^2)^{5/2}}
 &\le C(m)L\,\log_N^3,\\
 \sum_{|k|\ge\pi/\eps_N}\frac{(1+\log(2+|k|))^2}{(m^2+k^2)^{3/2}}
 &\le C(m)L\,\eps_N^2\log_N^2
 \end{split}
 \label{eq:kcount}
\end{equation}
(the first summand decays like $|k|^{-1}\log^2|k|$, so its band sum is
$O(\log^3)$; the second is a convergent tail of $|k|^{-3}\log^2|k|$).

\emph{Fourth chaos.} By Lemma~\ref{lem:threesums} and the decomposition of
$\E[X_4^2]$ into principal and alias parts,
\[
 \begin{aligned}
 \|X_4-Y_4\|_2^2
 &=4!\,(2Lt)^{-2}\,B_4+\mathrm{Al}_N-2\mathcal E_N^{\rm edge},\\
 B_4&:=\!\!\sum_{\substack{\kappa\in D_N^4\\ \Sigma\kappa=0}}\!\!
 \Bigl[\prod_i\hat c_N-2\prod_i\hat c_\times+\prod_i\hat c_\infty\Bigr]\\
 &\quad+\!\!\sum_{\substack{\kappa\in D_\infty^4\setminus D_N^4\\
 \Sigma\kappa=0}}\prod_i\hat c_\infty .
 \end{aligned}
\]
For the first sum, telescope each product difference over the four factors:
\[
 \Bigl|\prod\nolimits_i\hat c_a(\kappa_i)-\prod\nolimits_i
 \hat c_\infty(\kappa_i)\Bigr|
 \le\sum_{i=1}^4|\hat c_a-\hat c_\infty|(\kappa_i)
 \prod_{j\ne i}\hat c_{\mathrm{env}}(\kappa_j),
 \qquad a\in\{N,\times\},
\]
valid on $D_N^4$ by Lemma~\ref{lem:envelope}(i); with
$|\hat c_\times-\hat c_\infty|\le|\hat c_N-\hat c_\infty|$
(Lemma~\ref{lem:symdiff}) and the $i$-symmetry, and dominating the
constrained three-fold sums by $F_3$ (positive summands, larger index set),
\begin{gather*}
 |B_4|\ \le\ 12\,T_1+4\,T_2,\\
 T_1:=\sum_{\kappa_1\in D_N}|\hat c_N-\hat c_\infty|(\kappa_1)\,
 F_3(-\kappa_1),\qquad
 T_2:=\sum_{\kappa_1\in D_\infty\setminus D_N}
 \hat c_{\mathrm{env}}(\kappa_1)\,F_3(-\kappa_1).
\end{gather*}
By Lemma~\ref{lem:symdiff}, $|\hat c_N-\hat c_\infty|(\kappa)
\le\frac{\pi^6}{768}\,\eps_N^2k^4(m^2+|\kappa|^2)^{-2}$, so with
\eqref{eq:convdecay}, \eqref{eq:k0count} (at $s=3$) and \eqref{eq:kcount},
\begin{align*}
 T_1&\le C(m)(tL)^2\eps_N^2\sum_{\kappa\in D_N}
 \frac{k^4(1+\log(2+|\kappa|))^2}{(m^2+|\kappa|^2)^3}\\
 &\le C(m,t_0)\,(tL)^2\,t\,\eps_N^2
 \sum_{k\in\Gamma_N}\frac{k^4(1+\log(2+|k|))^2}{(m^2+k^2)^{5/2}}
 \ \le\ C\,(tL)^3\eps_N^2\log_N^3,
\end{align*}
and with \eqref{eq:k0count} at $s=2$:
\begin{align*}
 T_2&\le C(m)(tL)^2\sum_{\substack{\kappa\in D_\infty\\ |k|\ge\pi/\eps_N}}
 \frac{(1+\log(2+|\kappa|))^2}{(m^2+|\kappa|^2)^2}\\
 &\le C(m,t_0)(tL)^2\,t\sum_{|k|\ge\pi/\eps_N}
 \frac{(1+\log(2+|k|))^2}{(m^2+k^2)^{3/2}}
 \ \le\ C\,(tL)^3\eps_N^2\log_N^2 .
\end{align*}
With the prefactor, the alias bound \eqref{eq:alias}, and the separately
tracked edge bound \eqref{eq:Eedge-bound},
\begin{equation}
 \begin{split}
 \|X_4-Y_4\|_2^2
 &\le4!(2Lt)^{-2}|B_4|+\mathrm{Al}_N
       +2|\mathcal E_N^{\rm edge}|\\
 &\le C(m,t_0)\,tL\;\eps_N^2\,\log_N^3 .
 \end{split}
 \label{eq:X4rate}
\end{equation}
Here $t_0L\ge1$, $\eps_N\le1/2$, and $\log_N\ge1$ imply that both
$t\eps_N^3\log_N^2$ and the explicit
$9t\cth(2t_0)^3\eps_N^5/(256L^2)$ are bounded by the last line after
enlarging a constant depending only on $(m,t_0)$.

\emph{Second chaos.} Here the principal part is an exact square: by
Lemma~\ref{lem:isserlis} (as in Lemma~\ref{lem:threesums}, with net
prefactor $(2Lt)^0$), using $\hat c_\times(\kappa)\hat c_\times(-\kappa)
=\hat c_N(\kappa)\hat c_\infty(\kappa)$ and evenness of all symbols,
\[
 \|X_2-Y_2\|_2^2
 =2\sum_{\kappa\in D_N}\bigl(\hat c_N-\hat c_\infty\bigr)^2(\kappa)
 +2\sum_{\kappa\in D_\infty\setminus D_N}\hat c_\infty(\kappa)^2 ,
\]
\emph{exactly}: no alias term and no edge remainder survive. Indeed, on the edge
column the lattice alias pairing $k_1=k_2=k_e$ --- forced by the unique alias value
$-2\pi/\eps_N$, since $-k_e\notin\Gamma_N$ --- supplies precisely the principal
pairing that the asymmetry withholds, so
$\E[X_2^2]=2\sum_{\kappa\in D_N}\hat c_N(\kappa)^2$ with the alias \emph{replacing},
not augmenting, the missing term; and by Proposition~\ref{prop:laws}(iii)(a) the two
symmetrized cross configurations $(\pm\pi/\eps_N,\mp\pi/\eps_N)$, each of weight
$\chi^2=\tfrac12$, restore full weight, so
$\E[X_2Y_2]=2\sum_{\kappa\in D_N}\hat c_N\hat c_\infty$. With
$\E[Y_2^2]=2\sum_{D_\infty}\hat c_\infty^2$, the display follows by subtraction. The square term is bounded by
Lemma~\ref{lem:symdiff} (squared: $(\hat c_N-\hat c_\infty)^2\le
C\eps_N^4k^8(m^2+|\kappa|^2)^{-4}$) together with \eqref{eq:k0count} at
$s=4$, which gives $\sum_{k_0}\le C(m,t_0)t\,(m^2+k^2)^{-7/2}$; the
remaining $k$-sum $\sum_{\Gamma_N}k^8(m^2+k^2)^{-7/2}$ has summand
$\sim|k|$ and band value $O(L\eps_N^{-2})$, so
$2\sum_{D_N}(\hat c_N-\hat c_\infty)^2\le C(m,t_0)\,tL\,\eps_N^{2}$ --- the
same order as the missing-mode term, not subleading;
the missing modes give $\le C(m,t_0)tL\,\eps_N^2$ by \eqref{eq:k0count} ($s=2$)
followed by the log-free tail
$\sum_{|k|\ge\pi/\eps_N}(m^2+k^2)^{-3/2}\le C(m)L\eps_N^2$, which is the second line
of \eqref{eq:kcount} with the numerator $(1+\log(2+|k|))^2$ deleted --- the same
integral comparison, run without the logarithm. Hence
$\|X_2-Y_2\|_2^2\le C(m,t_0)\,tL\,\eps_N^2$.

\emph{Assembly.} By Definition~\ref{def:U},
\[
 \|\U_N-\U\|_2\le\lambda\Bigl[\|X_4-Y_4\|_2+6\bar d\,\|X_2-Y_2\|_2
 +6|d_N-d_\infty|\,\|Y_2\|_2+3|d_N^2-d_\infty^2|\,(2Lt)\Bigr],
\]
with $\bar d$ from Lemma~\ref{lem:recomb}, $\|Y_2\|_2^2=2S_2^{(\infty)}\le
2\Sigma_1\le C(m)tL$, and $|d_N-d_\infty|\le C\eps_N^2$,
$|d_N^2-d_\infty^2|\le2\bar d\,C\eps_N^2$. All four terms are
$\le C(m,\lambda,t_0)(1+t)(1+L)\eps_N\log_N^{3/2}$ (for the last,
$\eps_N^2\,t\le\eps_N(1+t)$ since $\eps_N\le1$), which is \eqref{eq:A4a}.
\end{proof}

\begin{remark}
All three mechanisms --- dispersion ($T_1$), missing modes ($T_2$), aliasing
--- contribute $O(\eps_N^2\log_N^{2\text{--}3})$ to the squared norm; none
dominates parametrically, and no exponent was given away. For A6 only
$\|\U_N-\U\|_{L^2}\to0$ is consumed; the explicit rate feeds the scale
selection in Theorem~\ref{thm:A5} and is recorded for any later quantitative
use. No numerical input is used anywhere in this section.
\end{remark}

\section{The renormalization dictionary (A4(b))}\label{sec:dict}

\begin{theorem}[the dictionary constant, with rate]\label{thm:dict}
Let $c^{\mathrm{ct}}_{L,N}:=\frac1{2L}\sum_{k\in\Gamma_N}\frac1{2\gamma(k)}$
(the $\Gamma_N$-truncated continuum vacuum constant). Then
\begin{equation}
 \Bigl|\ \delta_N-\frac1{2\pi}\log\frac4\pi\ \Bigr|
 \;\le\;C(m,L)\;\eps_N
 \qquad(\eps_N\le\tfrac12),
 \qquad
 \delta_N:=c_{N,L}-c^{\mathrm{ct}}_{L,N},
 \label{eq:dict}
\end{equation}
with $C(m,L)$ explicit from the proof
($C=\tfrac1{8L}+\tfrac{\pi C_0}{8L}+C'(m)$, $C_0<6$: the Riemann step
contributes $C_0\tfrac\pi2\Delta u$ with $\Delta u=\pi\eps_N/(2L)$, carried
through the prefactor $\tfrac1{2\pi}$ of Step 1, i.e.\
$\tfrac1{2\pi}\,C_0\,\tfrac\pi2\cdot\tfrac{\pi\eps_N}{2L}
=\tfrac{\pi C_0}{8L}\,\eps_N$).
\end{theorem}

\begin{proof}
\emph{Step 1: exact reduction.} With $k_j=\pi j/L$, $u_j:=\eps_Nk_j/2
=\pi j/(2r_N)$ and $\tilde\delta:=m\eps_N/2$:
$\gamma_N(k_j)=\tfrac2{\eps_N}(\tilde\delta^2+\sin^2u_j)^{1/2}$ and
$\gamma(k_j)=\tfrac2{\eps_N}(\tilde\delta^2+u_j^2)^{1/2}$ exactly, so
\[
 \delta_N=\frac1{4L}\sum_{j=-r_N}^{r_N-1}
 \Bigl(\frac1{\gamma_N}-\frac1\gamma\Bigr)(k_j)
 =\frac1{2\pi}\sum_{j=1}^{r_N-1}g_{\tilde\delta}(u_j)\,\Delta u
 +E_{\rm edge},
\]
where $\Delta u=\pi/(2r_N)$,
$g_\delta(u):=(\delta^2+\sin^2u)^{-1/2}-(\delta^2+u^2)^{-1/2}$
(the $j=0$ term vanishes; the terms $\pm j$, $1\le j\le r_N-1$, are equal and
were paired, producing the factor $2\cdot\frac1{4L}\cdot\frac{\eps_N}2
=\frac{\Delta u}{2\pi}$ per term). The single unpaired edge term $j=-r_N$ is
$E_{\rm edge}=\frac1{4L}\cdot\frac{\eps_N}2\,g_{\tilde\delta}(\pi/2)$, and
$0\le g_{\tilde\delta}(\pi/2)\le1$, so $0\le E_{\rm edge}\le\eps_N/(8L)$.

\emph{Step 2: uniform $C^1$ bound.} Write $g_\delta(u)
=\psi(\sin^2u)-\psi(u^2)$ with $\psi(a)=(\delta^2+a)^{-1/2}$. Then
\[
 g_\delta'(u)=\bigl[\psi'(\sin^2u)-\psi'(u^2)\bigr]\sin2u
 +\psi'(u^2)\,(\sin2u-2u).
\]
On $(0,\pi/2]$: $|\psi'(a)-\psi'(b)|\le\tfrac34(\min(a,b))^{-5/2}|a-b|$
(from $\psi''(a)=\tfrac34(\delta^2+a)^{-5/2}$ and the mean value theorem),
$u^2-\sin^2u\le u^4/3$, $\sin^2u\ge(2u/\pi)^2$, $|\sin2u|\le2u$,
$|\psi'(u^2)|\le\tfrac12u^{-3}$, $|\sin2u-2u|\le\tfrac{(2u)^3}6$; hence
\[
 |g_\delta'(u)|\le\tfrac34\Bigl(\tfrac\pi{2u}\Bigr)^5\cdot\tfrac{u^4}3
 \cdot2u+\tfrac12u^{-3}\cdot\tfrac{8u^3}6
 =\tfrac12\Bigl(\tfrac\pi2\Bigr)^5+\tfrac23=:C_0<6,
\]
uniformly in $u\in(0,\pi/2]$ and $\delta\ge0$ (and $g_\delta(0^+)=0$). Therefore the left-endpoint Riemann sum obeys
\[
 \Bigl|\sum_{j=1}^{r_N-1}g_{\tilde\delta}(u_j)\Delta u
 -\int_0^{\pi/2}g_{\tilde\delta}(u)\,du\Bigr|
 \le C_0\,\tfrac\pi2\,\Delta u\;\le\;C\,\eps_N .
\]

\emph{Step 3: removing $\tilde\delta$.} With
$\varphi(a):=a^{-1/2}-(\delta^2+a)^{-1/2}
=\int_0^{\delta^2}\tfrac12(s+a)^{-3/2}ds\ge0$:
$g_0(u)-g_\delta(u)=\varphi(\sin^2u)-\varphi(u^2)\in[0,\,
(u^2-\sin^2u)\max_{[\sin^2u,u^2]}|\varphi'|\,]$, and
$|\varphi'(a)|=\tfrac12|a^{-3/2}-(\delta^2+a)^{-3/2}|
\le\min\bigl(\tfrac12a^{-3/2},\ \tfrac34\delta^2a^{-5/2}\bigr)$. For
$u\le\delta$ use the first bound:
$g_0-g_\delta\le\tfrac{u^4}3\cdot\tfrac12(2u/\pi)^{-3}\le Cu$; for
$u\ge\delta$ use the second:
$g_0-g_\delta\le\tfrac{u^4}3\cdot\tfrac34\delta^2(2u/\pi)^{-5}
\le C\delta^2/u$. Hence
\[
 0\le\int_0^{\pi/2}\bigl(g_0-g_{\tilde\delta}\bigr)
 \le C\tilde\delta^2+C\tilde\delta^2\log\frac\pi{2\tilde\delta}
 \;\le\;C'\eps_N^2\log(1/\eps_N)\qquad(\eps_N\le\tfrac12).
\]

\emph{Step 4: the integral.} $\int_0^{\pi/2}(\csc u-u^{-1})du
=\bigl[\log\tan(u/2)-\log u\bigr]_0^{\pi/2}
=-\log\tfrac\pi2-\log\tfrac12=\log\tfrac4\pi$, since
$\tan(u/2)/u\to\tfrac12$ as $u\downarrow0$ and $\tan(\pi/4)=1$. Combining
Steps 1--4 proves \eqref{eq:dict}. (V0: the numeric value
$0.0384461803$ was confirmed to $3.6\times10^{-10}$; the observed
approach $\approx\eps_N^2$ is faster than the proved $O(\eps_N)$ --- the
proof's Riemann-sum step is first-order because $g_0'$ does not vanish at
$\pi/2$; we do not need to sharpen it.)
\end{proof}

\begin{corollary}[the convention lock: no counterterm survives]
\label{cor:lock}
Let ${:}\;{:}^{\rm tr}$ denote the continuum Wick convention defined by the
$\Gamma$-truncation limit (subtract $c^{\mathrm{ct}}_{L,K}$ before
$K\to\infty$). Then the continuum interaction $\U$ of Definition~\ref{def:U}
\emph{is} the truncation-convention quartic,
\[
 \U=\lambda\int_0^t\!\!\int_{\T_L}{:}\eta^4{:}^{\rm tr},
\]
with \emph{no} quadratic and no constant residue: the gate's continuum
polynomial is plain $\varphi^4$. Equivalently, the lattice-vacuum convention
of the model sheet and the continuum truncation convention \emph{agree in the
limit}: although the two Wick constants differ by
$\delta_\infty=\tfrac1{2\pi}\log\tfrac4\pi>0$ (Theorem~\ref{thm:dict}), the
two fields' equal-point variances differ by the same $\delta_\infty$, and the
two offsets cancel.
\end{corollary}

\begin{proof}
By Definition~\ref{def:U} and Lemma~\ref{lem:recomb} (recombination in the
direction $v\to c$), at spatial truncation $K$,
\[
 \tfrac1\lambda\U^{(K)}
 =Y^{(K)}_4+6d_\infty Y^{(K)}_2+3d_\infty^2\,\mathrm{Vol}
 =\int_0^t\!\!\int_{\T_L}{:}\eta_K^4{:}_{c}\,,
 \qquad c=v^{(K)}_\infty-d_\infty ,
\]
where $v^{(K)}_\infty(t)=\tfrac1{2L}\sum_{|k|\le K}
\cth(t\gamma(k)/2)/(2\gamma(k))$ is the truncated continuum field's own
variance (the chaoses $Y^{(K)}_j$ are ordered by it,
Proposition~\ref{prop:Yexists}). Now
\[
 v^{(K)}_\infty-\frac1{2L}\!\sum_{|k|\le K}\!\frac{\cth(t\gamma/2)-1}{2\gamma}
 =\frac1{2L}\sum_{|k|\le K}\frac1{2\gamma(k)}
 =c^{\mathrm{ct}}_{L,K},
\]
so, with
$d^{(K)}_\infty:=\tfrac1{2L}\sum_{|k|\le K}
[\gamma(k)(e^{t\gamma(k)}-1)]^{-1}$ the $K$-truncated continuum thermal
constant, $c=c^{\mathrm{ct}}_{L,K}+(d^{(K)}_\infty-d_\infty)$ with
\[
 0\le d_\infty-d^{(K)}_\infty
 =\frac1{2L}\sum_{|k|>K}\frac1{\gamma(k)(e^{t\gamma(k)}-1)}
 \le C(m,L,t_0)\,e^{-t_0K}
\]
(using $\gamma(k)\ge|k|$), and letting $K\to\infty$ gives the display.

The second formulation is the same computation read on the lattice side. By
\eqref{eq:dN}, $d_N=v_N-c_{N,L}$ with $v_N$ the lattice field's own variance,
so
\[
 v_N-v^{(K)}_\infty
 =\bigl(c_{N,L}-c^{\mathrm{ct}}_{L,K}\bigr)+\bigl(d_N-d^{(K)}_\infty\bigr)
 =\delta_N+\bigl(d_N-d^{(K)}_\infty\bigr)=\delta_N+O(\eps_N^2)
 \ \longrightarrow\ \delta_\infty
\]
with the $K$-truncation taken to be $\Gamma_N$ itself (so that
$c^{\mathrm{ct}}_{L,K}=c^{\mathrm{ct}}_{L,N}$ and the displayed identity is
exact), by Theorem~\ref{thm:dict} and Lemma~\ref{lem:recomb}. Taking instead
the symmetric band $\{|k|\le\pi/\eps_N\}$ adds the single edge mode
$k=+\pi/\eps_N$, worth $1/(4L\gamma(\pi/\eps_N))=\eps_N/(4\pi L)+O(\eps_N^3)$,
so that the rate becomes $\delta_N+O(\eps_N)$; the limit $\delta_\infty$,
which is all the corollary consumes, is unchanged. The
excess tadpole of the lattice field is thus exactly the excess Wick
subtraction it receives, which is why $\U_N$ and $\U$ --- each
\emph{completely} Wick-ordered by its own field's vacuum constant --- are the
matched pair compared in Theorem~\ref{thm:A4a}, with no counterterm between
them.
\end{proof}

\begin{remark}[what $\delta_\infty$ is, and is not]\label{rem:deltarole}
$\delta_\infty=\tfrac1{2\pi}\log\tfrac4\pi$ is a genuine, correctly computed
constant: it is the offset between the two \emph{tadpoles}
$c_{N,L}$ and $c^{\mathrm{ct}}_{L,N}$ (Theorem~\ref{thm:dict}, confirmed
numerically in V0 to $3.6\times10^{-10}$). It is \emph{not} a finite mass
renormalization of the gate's continuum polynomial. An earlier version of
this corollary asserted the residue
$\varphi^4-6\delta_\infty\varphi^2+3\delta_\infty^2$; that is false, because
it compares the two Wick \emph{constants} while holding the field fixed,
whereas the lattice and continuum fields have variances differing by the same
$\delta_\infty$. The error was found in the round-2 adversarial pass
(independently by two reviewers, confirmed by four verifiers) and is recorded
in \texttt{strategy.md}; the correction propagates to V3 \S1, whose gate
endpoint is $(a,b)=(0,0)$.
\end{remark}

\section{Mean-shift coefficients (A4(c))}\label{sec:A4c}

Recall from V1 (Lemma 5.2, \texttt{lem:interp}): for the Weyl symbol
$\xi=\eps_Mq+ip$ at coarse scale $M$, the mean-shift representation at scale
$N=M+r$ shifts the interaction by the harmonic jump interpolant, per mode
$H_k(s)=A_k\cosh(\gamma_N(k)s)+B_k\sh(\gamma_N(k)s)$ with $A_k=-\rho_k/2$,
$B_k=\tfrac{\rho_k}2\cth(t\gamma_N(k)/2)$; the refined momentum symbol is
$\widehat{p_N}(k)=2^{r/2}P_r(\eps_Mk)\,\widehat{p_M}(k)$ on $\Gamma_N$
(coarse hat $2\pi/\eps_M$-periodic; free-rate document \S1), with the $D4$
envelope $|P_r(\ell)|^2\le C_D(1+|\ell|)^{-2-\eta}$, $\eta=2-\log_23$,
uniformly in $r\le\infty$ (free-rate Lemma 2.1, eq.\ (2.4)).

\begin{lemma}[closed form of the interpolant]\label{lem:closed}
For every $\omega>0$ and jump $\rho$, the interpolant of V1, Lemma 5.2
equals
\begin{equation}
 H(s)=\frac\rho2\,\frac{\sh\bigl(\omega(s-\tfrac t2)\bigr)}
 {\sh(\omega t/2)},\qquad s\in[0,t],
 \label{eq:closedform}
\end{equation}
and consequently $|H(s)|\le|\rho|/2$ for all $s\in[0,t]$, and
$|\partial_\omega H(s)|\le\tfrac{|\rho|}2\,t\,\cth(\omega t/2)$.
\end{lemma}

\begin{proof}
The hyperbolic addition formula gives
\[
 \sh\bigl(\omega s-\tfrac{\omega t}2\bigr)
 =\sh(\omega s)\ch\bigl(\tfrac{\omega t}2\bigr)
 -\ch(\omega s)\sh\bigl(\tfrac{\omega t}2\bigr);
\]
dividing by $\sh(\omega t/2)$ and multiplying by $\rho/2$ yields exactly
\[
 -\tfrac\rho2\ch(\omega s)
 +\tfrac\rho2\cth\bigl(\tfrac{\omega t}2\bigr)\sh(\omega s)
 =A\ch(\omega s)+B\sh(\omega s).
\] The bound $|H|\le|\rho|/2$ holds because $|\sh|$ is even and
increasing and $|s-t/2|\le t/2$. For the $\omega$-derivative, with
$u:=s-t/2$,
\[
 \partial_\omega H(s)=\frac\rho2\,
 \frac{u\ch(\omega u)\sh(\tfrac{\omega t}2)
 -\tfrac t2\ch(\tfrac{\omega t}2)\sh(\omega u)}{\sh^2(\omega t/2)},
\]
and each of the two terms is bounded in modulus by
$\tfrac{|\rho|}2\cdot\tfrac t2\cth(\omega t/2)$ (use $|u|\le t/2$,
$\ch(\omega u)\le\ch(\omega t/2)$, $|\sh(\omega u)|\le\sh(\omega t/2)$).
\end{proof}

\begin{definition}[shift fields]\label{def:shift}
In the original field coordinates (lattice transform
$\widehat f(k)=\eps_N^{1/2}\sum_xf(x)e^{-ikx}$, inverse
$f(x)=\tfrac{\eps_N^{1/2}}{2L}\sum_k\widehat f(k)e^{ikx}$), set
$\phi_\gamma(s):=\sh(\gamma(s-t/2))/\sh(\gamma t/2)$ and, using
$\eps_N^{1/2}2^{r/2}=\eps_M^{1/2}$,
\begin{align*}
 h^{(N)}(s,x)&:=\frac1{2L}\sum_{k\in\Gamma_N}a^{(N)}_k(s)\,e^{ikx},&
 a^{(N)}_k(s)&:=\frac{\eps_M^{1/2}}2\,P_r(\eps_Mk)\,\widehat{p_M}(k)\,
 \phi_{\gamma_N(k)}(s),\\
 h^{(\infty)}(s,x)&:=\frac1{2L}\sum_{k\in\Gamma_\infty}a^{(\infty)}_k(s)
 \,e^{ikx},&
 a^{(\infty)}_k(s)&:=\frac{\eps_M^{1/2}}2\,P_\infty(\eps_Mk)\,
 \widehat{p_M}(k)\,\phi_{\gamma(k)}(s).
\end{align*}
By Lemma~\ref{lem:closed}, $h^{(N)}$ is precisely V1's interpolant path read
in the original coordinates (same per-mode harmonic profile, boundary data
$h^{(N)}(0)=-p_N/2$, $h^{(N)}(t)=+p_N/2$, periodic derivative; V1's mode
relation is linear and inverts, as noted there); $h^{(\infty)}$ is its
continuum counterpart with the free-rate multiplier $P_\infty$.
\end{definition}

\begin{remarknn}[the $D4$ filter annihilates the edge column; reality of $h^{(N)}$]
The mask satisfies $m_0(\pi)=0$: by the model sheet (2.3),
$m_0(\vartheta)=\bigl(\tfrac{1+e^{-i\vartheta}}2\bigr)^2D(\vartheta)$, and
$1+e^{-i\pi}=0$. Since $\eps_Mk_e=-\pi\eps_M/\eps_N=-\pi2^{\,r}$, the product
$P_r(\eps_Mk_e)=\prod_{j=1}^rm_0(2^{-j}\eps_Mk_e)$ contains the factor
$m_0(-\pi)=0$, so
\[
 P_r(\pm\pi2^{\,r})=0\qquad(r\ge1),
\]
and therefore $a^{(N)}_{k_e}\equiv0$ for every $r\ge1$: \emph{the coarse data never
excite the edge column.} (For $r=0$, $P_0=1$ and $\widehat{p_M}(-\pi/\eps_M)$ is real
because $p_M$ is.) Two consequences are used below.
\emph{(a) $h^{(N)}$ and $h^{(\infty)}$ are real.} On $\Gamma_\infty$, and on
$\Gamma_N\setminus\{k_e\}$, the Hermitian pairing $a_{-k}=\overline{a_k}$ closes
($p_M$ real, $P_r$ real-masked, $\gamma_N$ and $\gamma$ even), exactly as for
$\eta$; the one unpaired mode $k_e$ carries the real function
$e^{ik_ex}=(-1)^{x/\eps_N}$ on $\Lambda_N$ and the coefficient $a^{(N)}_{k_e}$, which
vanishes for $r\ge1$ and is real for $r=0$. So no analogue of the defect of
Definition~\ref{def:coupling} arises here. This is what licenses the shift clause of
Lemma~\ref{lem:lower}, whose pointwise identity is between \emph{real} numbers.
\emph{(b) Where the edge remainder $R_N$ drops out, and where it does not.} $R_N$ is
supported on the edge column, so it is annihilated by pairing against a coefficient with
no edge mode --- but \emph{only where the mixed kernel enters linearly}. That is the case
for V4's phase functional $\Phi_N(J_q)$, which is paired directly against $a^{(N)}$ and
therefore loses its edge contribution for $r\ge1$. It is \emph{not} the case in
Proposition~\ref{prop:shifted}, where the kernel is raised to the chaos order $j$. Write
$R_N(z,z')=\sigma(x)\Psi(s-s',y)$ with $\sigma(x)=e^{ik_ex}=(-1)^{x/\eps_N}$ on
$\Lambda_N$, so that $\sigma^2\equiv1$. Then, e.g.\ for $j=3$,
\[
 (C_\times+R_N)^3-C_\times^3
 =3C_\times^2\sigma\Psi+3C_\times\Psi^2+\sigma\Psi^3 ,
\]
whose middle term carries \emph{no} $\sigma$ and is supported on the whole
three-fold sumset: it is not edge-supported, and no vanishing of $a^{(N)}_{k_e}$
annihilates it.
Only the kernel power $1$ is edge-supported, and there $g_N=(h^{(N)})^3$, whose
edge Fourier coefficient is a \emph{convolution},
$\sum_{k_1+k_2+k_3\equiv k_e}\prod_ia^{(N)}_{k_i}$, taken over configurations with
no factor on the edge column and hence generically nonzero. In all three chaoses of
Proposition~\ref{prop:shifted} the remainder is therefore \emph{bounded, not
eliminated}; the bound is in that proposition's proof and covers $j\in\{1,2,3\}$
uniformly.
\end{remarknn}

\begin{lemma}[shift-field bounds; symbol convergence]
\label{lem:H}
With $\bar P_M:=\max_{k\in\Gamma_M}|\widehat{p_M}(k)|$, set
\[
 \Sigma_M(p):=\frac{C_D^{1/2}\,\eps_M^{1/2}}{4L}
 \Bigl(1+\frac{4L}{\pi\eta\,\eps_M}\Bigr)\bar P_M .
\]
Then:
\emph{(i)} $\sup_N\sup_{s\in[0,t],\,x}|h^{(N)}(s,x)|\le\Sigma_M(p)$ and
$\sup_{s,x}|h^{(\infty)}(s,x)|\le\Sigma_M(p)$;
\emph{(ii)} for every $j\ge1$,
$\eps_N\sum_{x}\int_0^t|h^{(N)}|^{\,j}\,ds\le2Lt\,\Sigma_M(p)^j$,
and likewise for $h^{(\infty)}$ with $\int_{\T_L}dx$;
\emph{(iii)} with $\theta_H:=(1+\eta)/2$,
\[
 \sup_{s\in[0,t]}\ \Bigl[\frac1{2L}\sum_{k\in\Gamma_N}
 \bigl|a^{(N)}_k(s)-a^{(\infty)}_k(s)\bigr|^2
 +\frac1{2L}\sum_{|k|\ge\pi/\eps_N}\bigl|a^{(\infty)}_k(s)\bigr|^2\Bigr]
 \ \le\ C_H\;2^{-\theta_H\,r},
\]
where $C_H=C(m,\eta)\,\bigl(1+t\,\cth(t_0m/2)\bigr)^2\,
\eps_M^{-4}(1+\eps_M)^{5}\,\bar P_M^{\,2}\,/\,L\cdot
\bigl(1+\tfrac L{\eps_M}\bigr)$ is independent of $N$ and $s$.
\end{lemma}

\begin{proof}
(i) By Lemma~\ref{lem:closed}, $|\phi_\gamma|\le1$, so
\[
 \begin{split}
 |h^{(N)}(s,x)|&\le\frac{\eps_M^{1/2}}{4L}\,\bar P_M
 \sum_{k\in\Gamma_N}|P_r(\eps_Mk)|\\
 &\le\frac{\eps_M^{1/2}}{4L}\,\bar P_M\,C_D^{1/2}
 \Bigl(1+\frac{2L}\pi\int_0^\infty\!\!(1+\eps_Mu)^{-1-\eta/2}du\Bigr)
 =\Sigma_M(p),
 \end{split}
\]
using the $k$-spacing $\pi/L$ and $\int_0^\infty(1+\eps_Mu)^{-1-\eta/2}du
=\tfrac2{\eta\eps_M}$; the same count over $\Gamma_\infty$ covers
$h^{(\infty)}$. (ii) is immediate from (i) ($\eps_N\cdot2r_N\cdot t=2Lt$).

(iii) Write $\ell:=\eps_Mk$ and split
\[
 a^{(N)}_k-a^{(\infty)}_k=\frac{\eps_M^{1/2}}2\widehat{p_M}(k)
 \Bigl[(P_r-P_\infty)(\ell)\,\phi_{\gamma_N(k)}
 +P_\infty(\ell)\bigl(\phi_{\gamma_N(k)}-\phi_{\gamma(k)}\bigr)\Bigr].
\]

\emph{Filter tail.} $P_\infty(\ell)=P_r(\ell)\prod_{j>r}m_0(2^{-j}\ell)$,
$|m_0|\le1$ (QMF), and $|1-\prod_jz_j|\le\sum_j|1-z_j|$ for $|z_j|\le1$
(telescoping), while $|1-m_0(x)|=|\int_0^xm_0'|\le L_{m_0}|x|$ with
$L_{m_0}:=\sup|m_0'|\le\tfrac1{\sqrt2}\sum_n n|h_n|=\tfrac{3+\sqrt3}4$ for
the $D4$ mask. The first two members of \eqref{eq:filtertail} below hold for all
$\ell$; the last uses the envelope of Lemma 2.1 of the free-rate document, which is
proved only on the band, so \eqref{eq:filtertail} is asserted for
$|\ell|\le\pi2^r$ --- which is the only range in which it is used here and in V4. (It
genuinely fails outside: $P_r$ is $2^{r+1}\pi$-periodic with $P_r(2^{r+1}\pi)=1$, while
$m_0(\pi)=0$ forces $P_\infty(2^{r+1}\pi)=0$, so the left side equals $1$ there whereas
the right side tends to $0$.) Hence, for $|\ell|\le\pi2^r$,
\begin{equation}
 |P_r-P_\infty|(\ell)\le|P_r(\ell)|\,L_{m_0}|\ell|\sum_{j>r}2^{-j}
 =L_{m_0}\,2^{-r}\,|\ell|\,|P_r(\ell)|
 \le C_D^{1/2}L_{m_0}2^{-r}|\ell|(1+|\ell|)^{-1-\eta/2}.
 \label{eq:filtertail}
\end{equation}
Its contribution to the mode sum is, with $|\phi|\le1$ and the band
restriction $|\ell|\le\pi2^r$,
\[
 \frac{\eps_M\bar P_M^2}{8L}\,C_DL_{m_0}^2\,4^{-r}
 \sum_{k\in\Gamma_N}\ell^2(1+\ell)^{-2-\eta}
 \le C(\eta)\,C_DL_{m_0}^2\,\frac{\bar P_M^2}{4\pi}\,
 4^{-r}\,(\pi2^r)^{1-\eta}
 \le C\,\bar P_M^2\,2^{-r(1+\eta)},
\]
by the $k$-count
\[
 \sum_{k\in\Gamma_N,\ \eps_M|k|\le X}G(\eps_M|k|)\ \le\
 \tfrac{2L}{\pi\eps_M}\int_0^XG+2\sup_{[0,X]}G
\]
valid for any $G\ge0$ that is \emph{unimodal} on $[0,X]$ (nondecreasing then
nonincreasing) --- which both integrands used here are --- and applied to
$G(\ell)=\ell^2(1+\ell)^{-2-\eta}$ with $X=\pi2^r$
($\int_0^X\ell^{-\eta}d\ell=X^{1-\eta}/(1-\eta)$). The correction term must be
$\sup_{[0,X]}G$, not $G(0)$: the left-endpoint comparison
$\sum_jG(j\Delta)\le\Delta^{-1}\!\int G+G(0)$ is valid only for \emph{nonincreasing}
$G$, whereas the integrands here vanish at $0$ and increase over most of the range.
The correction is harmless: it contributes
$C\bar P_M^2\,4^{-r}\sup G\le C\bar P_M^2\,4^{-r}$, and $4^{-r}<2^{-r(1+\eta)}$.

\emph{Dispersion.} By Lemma~\ref{lem:closed} and the mean value theorem in
$\omega\in[\gamma_N(k),\gamma(k)]\subset[m,\infty)$, together with free-rate
Lemma 3.1 (eq.\ (3.1)),
\[
 |\phi_{\gamma_N(k)}(s)-\phi_{\gamma(k)}(s)|
 \le t\,\cth(t_0m/2)\,\frac{\eps_N^2k^4}{24m}
 \quad\text{and always}\quad
 |\phi_{\gamma_N}-\phi_\gamma|\le2 .
\]
Split at $\ell_r:=2^{r/2}$. For $\ell\le\ell_r$ use the first bound, noting
$\eps_N^2k^4=4^{-r}\eps_M^2k^4=4^{-r}\ell^4\eps_M^{-2}$:
\begin{multline*}
 \frac{\eps_M\bar P_M^2}{8L}\,C_D
 \Bigl(\frac{t\cth(\tfrac{t_0m}2)}{24m\,\eps_M^{2}}\Bigr)^{2}16^{-r}
 \sum_{k\in\Gamma_N:\ \ell\le\ell_r}\ell^8(1+\ell)^{-2-\eta}\\
 \le\ C(m,\eta)\,\frac{(1+t)^2\bar P_M^2}{\eps_M^{4}}
 \Bigl(1+\frac L{\eps_M}\Bigr)\frac1L\ 16^{-r}\,2^{\frac r2(7-\eta)},
\end{multline*}
by the $k$-count with $\int_0^{X}\ell^{6-\eta}d\ell=X^{7-\eta}/(7-\eta)$;
the exponent is $-4r+\tfrac r2(7-\eta)=-\tfrac r2(1+\eta)$, i.e.\ this term
is $O(2^{-r(1+\eta)/2})$. For $\ell>\ell_r$ use
$|\Delta\phi|\le2$ and the envelope:
$\tfrac{\eps_M\bar P_M^2}{2L}C_D\sum_{\ell>\ell_r}(1+\ell)^{-2-\eta}
\le C(\eta)\bar P_M^2\,2^{-r(1+\eta)/2}$
($\int_X^\infty(1+\ell)^{-2-\eta}\le X^{-1-\eta}$, and
$2^{-\frac r2(1+\eta)}$ majorizes it).

\emph{Missing modes.} $|k|\ge\pi/\eps_N$ means $\ell>\pi2^r$:
$\tfrac{\eps_M\bar P_M^2}{8L}C_D\sum_{\ell>\pi2^r}(1+\ell)^{-2-\eta}
\le C\bar P_M^2\,2^{-r(1+\eta)}$.

Collecting the three contributions (each uniform in $s$) and absorbing the
grid-count prefactors into $C_H$ proves (iii).
\end{proof}

\begin{lemma}[folded, exact-output, and high-output power symbols]
\label{lem:symbolsup}
Put $B_N:=\Gamma_N$, $\mathfrak R_N:=\pi/\eps_N$, and let $[q]_N$ denote
the unique representative of $q$ modulo $2\mathfrak R_N$ in the half-open band
$B_N=[-\mathfrak R_N,\mathfrak R_N)\cap(\pi/L)\Z$.  For $j\in\{1,2,3\}$ define
\begin{align}
 K_{N,j}^{\rm fold}(p_0,p)
 &:=(2Lt)^{-(j-1)}
 \sum_{\substack{\kappa_i\in D_N\\
       \sum_i k_{0,i}=p_0,\ [\sum_i k_i]_N=p}}
       \prod_i\hat c_N(\kappa_i),                                      \label{eq:Kfold}\\
 K_{\times,j}(q_0,q)
 &:=(2Lt)^{-(j-1)}
 \sum_{\substack{\kappa_i\in D_N\\\sum_i\kappa_i=(q_0,q)}}
       \prod_i\hat c_\times(\kappa_i),                                \label{eq:Kcross}\\
 K_{\infty,j}(q_0,q)
 &:=(2Lt)^{-(j-1)}
 \sum_{\substack{\kappa_i\in D_\infty\\\sum_i\kappa_i=(q_0,q)}}
       \prod_i\hat c_\infty(\kappa_i).                                \label{eq:Kinf}
\end{align}
An empty sum is zero.  Thus \eqref{eq:Kfold} is a multiplier on the lattice
fundamental band, whereas \eqref{eq:Kcross}--\eqref{eq:Kinf} retain the exact
continuum output.  For $\eps_N\le1/2$ and $t\ge t_0$:
\begin{align}
 \sup K_{N,j}^{\rm fold}+\sup K_{\times,j}+\sup K_{\infty,j}
 &\le B_j(m,t_0),                                                       \label{eq:Ksup}\\
 \sup_{p_0,\,p\in B_N}
 |K_{N,j}^{\rm fold}(p_0,p)-K_{\infty,j}(p_0,p)|
 &\le B_j'(m,t_0)\eps_N^2\log_N^2,                                    \label{eq:Kfolddiff}\\
 \sup_{p_0,\,p\in B_N}
 |K_{\times,j}(p_0,p)-K_{\infty,j}(p_0,p)|
 &\le B_j'(m,t_0)\eps_N^2\log_N^2,                                    \label{eq:Kcrossdiff}\\
 \sup_{q_0,\,|q|\ge\mathfrak R_N}
 [K_{\times,j}(q_0,q)+K_{\infty,j}(q_0,q)]
 &\le B_j''(m,t_0)\eps_N^2\log_N^{j-1}.                               \label{eq:Khigh}
\end{align}
All suprema are over the displayed temporal grid and spatial domains.
\end{lemma}

\begin{proof}
Let $\Phi_0:=\mathbf1_{\{0\}}$, $\Phi_1:=\hat c_{\rm env}$, and
$\Phi_2:=F_2$.  For $j=1$, \eqref{eq:Ksup} follows from
$\hat c_{\rm env}\le\pi^2/(4m^2)$.  For $j=2,3$, every exact-output sum is
bounded by $F_j$ (where $F_2,F_3$ are those of
Lemma~\ref{lem:envelope}).  The folded sum is the sum of the exact outputs
$(p_0,p+2n\mathfrak R_N)$ that lie in the $j$-fold band; there are at most
$j+1$ such integers $n$.  Equation \eqref{eq:convdecay} and the normalization
$(2Lt)^{-(j-1)}$ therefore prove \eqref{eq:Ksup}.

If $|q|\ge\mathfrak R_N$, then for $j=2,3$
\[
 K_{\times,j}(q_0,q)+K_{\infty,j}(q_0,q)
 \le2(2Lt)^{-(j-1)}F_j(q_0,q)
 \le C(m,t_0)\eps_N^2\log_N^{j-1}
\]
by \eqref{eq:convdecay}; for $j=1$, each nonzero summand is at most
$C(m^2+q^2)^{-1}\le C\eps_N^2$.  This proves \eqref{eq:Khigh}.

It remains to work on the fundamental band.  First restrict
$K_{\infty,j}$ to tuples in $D_N$.  Telescoping the $j$ factors and using
$|\hat c_\times-\hat c_\infty|\le|\hat c_N-\hat c_\infty|$ gives, for
$a=N$ or $a=\times$,
\begin{equation}
 \left|K_{a,j}^{\rm exact}(p)-K_{\infty,j}^{D_N}(p)\right|
 \le j(2Lt)^{-(j-1)}
 \sum_{u\in D_N}|\hat c_N-\hat c_\infty|(u)\Phi_{j-1}(p-u).
 \label{eq:Ktelescope}
\end{equation}
Here $K_{N,j}^{\rm exact}$ denotes the $n=0$ part of \eqref{eq:Kfold}.
Lemma~\ref{lem:symdiff} yields the flat bound
\begin{equation}
 |\hat c_N-\hat c_\infty|(u)
 \le\frac{\pi^6}{768}\eps_N^2
       \frac{u_{\rm sp}^4}{(m^2+|u|^2)^2}
 \le\frac{\pi^6}{768}\eps_N^2 .                                      \label{eq:symdiffflat}
\end{equation}
For $j=1$, \eqref{eq:Ktelescope} is therefore $O(\eps_N^2)$.  For
$j=2$, summing $\Phi_1$ over the doubled band gives
$C(m,t_0)tL\log_N$; for $j=3$, summing $\Phi_2$ gives
$C(m,t_0)(tL)^2\log_N^2$.  After the normalization in
\eqref{eq:Ktelescope},
\begin{equation}
 \sup_{p_0,p\in B_N}
 |K_{a,j}^{\rm exact}(p)-K_{\infty,j}^{D_N}(p)|
 \le C(m,t_0)\eps_N^2\log_N^{j-1}.                                  \label{eq:Kexactdiff}
\end{equation}

For completeness, the missing-input estimate used next is
\begin{equation}
 (2Lt)^{-(j-1)}
 \sum_{\substack{u\in D_\infty\\|u_{\rm sp}|\ge\mathfrak R_N}}
 \hat c_{\rm env}(u)\Phi_{j-1}(p-u)
 \le C(m,t_0)\eps_N^2\log_N^{j-1},qquad p_{\rm sp}\in B_N.           \label{eq:Kmissing}
\end{equation}
For $j=1$ this is interpreted as the single possible input and follows from
$\hat c_{\rm env}(p)\le C\eps_N^2$.  For $j=2,3$, sum the temporal
variable first.  With
$a=(m^2+u_{\rm sp}^2)^{1/2}$ and
$b=(m^2+(p_{\rm sp}-u_{\rm sp})^2)^{1/2}$, the elementary bound
\[
 \sum_{u_0\in(2\pi/t)\Z}
 \frac1{u_0^2+a^2}\frac1{(p_0-u_0)^2+b^2}
 \le\frac1{a^2}\sum_{u_0}\frac1{(p_0-u_0)^2+b^2}
 \le C(m,t_0)\frac{t}{a^2b}
\]
is uniform in $p_0$.  The spatial sum outside the band is split into
$\mathfrak R_N\le|u_{\rm sp}|\le2\mathfrak R_N$ and
$|u_{\rm sp}|>2\mathfrak R_N$.  On the first part it is bounded by
\[
 C t\mathfrak R_N^{-2}
 \sum_{0\le n\le C L\mathfrak R_N}
       \frac1{(m^2+(n\pi/L)^2)^{1/2}}
 \le C(m,t_0)tL\mathfrak R_N^{-2}\log_N;
\]
on the second, $|p_{\rm sp}-u_{\rm sp}|\ge|u_{\rm sp}|/2$ and the
summand is $O(t|u_{\rm sp}|^{-3})$, whose tail is
$O(tL\mathfrak R_N^{-2})$.  This proves \eqref{eq:Kmissing} for $j=2$.
For $j=3$ no appeal to a ``same split'' is needed.  First sum $F_2$ over
its temporal output.  Since
\[
 \sum_{k_0\in(2\pi/t)\Z}\hat c_{\rm env}(k_0,k)
 \le C(m,t_0)\frac{t}{(m^2+k^2)^{1/2}},
\]
Tonelli and the definition of $F_2$ give, with $d\in\Gamma_\infty$,
\begin{align}
 \sum_{q_0}F_2(q_0,d)
 &=\sum_{k\in\Gamma_\infty}
   \left(\sum_{k_0}\hat c_{\rm env}(k_0,k)\right)
   \left(\sum_{\ell_0}\hat c_{\rm env}(\ell_0,d-k)\right)\\
 &\le C(m,t_0)t^2
   \sum_{k\in\Gamma_\infty}
   \frac1{(m^2+k^2)^{1/2}(m^2+(d-k)^2)^{1/2}}\\
 &\le C(m,t_0)t^2L\,
   \frac{1+\log(2+|d|/m)}{(m^2+d^2)^{1/2}}.                       \label{eq:F2-temporal-sum}
\end{align}
For the last line split the $k$-sum into
$|k|\le|d|/2$, $|d-k|\le|d|/2$, and their complement.  On the first
region the second denominator is at least
$\frac12(m^2+d^2)^{1/2}$ and
$\sum_{|k|\le|d|/2}(m^2+k^2)^{-1/2}\le
C(m)L[1+\log(2+|d|/m)]$; the second region is identical after
$k\mapsto d-k$; on the complement both $|k|$ and $|d-k|$ are at least
$|d|/2$, and splitting further at $|k|=2|d|$ leaves a bounded middle
part plus a tail $\sum_{|k|>2|d|}k^{-2}\le CL/(m+|d|)$.  This proves
\eqref{eq:F2-temporal-sum}.

Now use $\hat c_{\rm env}(u_0,u_{\rm sp})\le
C(m^2+u_{\rm sp}^2)^{-1}$ and put $d=p_{\rm sp}-u_{\rm sp}$.  Translation
of the temporal grid and \eqref{eq:F2-temporal-sum} imply that the left side of
\eqref{eq:Kmissing} for $j=3$ is at most
\[
 \frac{C}{L}\sum_{|u_{\rm sp}|\ge\mathfrak R_N}
 \frac{1+\log(2+|p_{\rm sp}-u_{\rm sp}|/m)}
 {(m^2+u_{\rm sp}^2)(m^2+(p_{\rm sp}-u_{\rm sp})^2)^{1/2}}.
\]
On $\mathfrak R_N\le|u_{\rm sp}|\le2\mathfrak R_N$, the first
denominator is at least $c\mathfrak R_N^2$ and the remaining spatial sum
is at most $CL[1+\log(2+\mathfrak R_N/m)]^2$.  On
$|u_{\rm sp}|>2\mathfrak R_N$, $|p_{\rm sp}-u_{\rm sp}|\ge
|u_{\rm sp}|/2$, so the summand is at most
$C[1+\log(2+|u_{\rm sp}|/m)]|u_{\rm sp}|^{-3}$ and its tail is
$CL\mathfrak R_N^{-2}[1+\log(2+\mathfrak R_N/m)]$.  Since
$\mathfrak R_N^{-2}=\pi^{-2}\eps_N^2$, this is precisely
$C(m,t_0)\eps_N^2\log_N^2$.  Hence \eqref{eq:Kmissing} holds for $j=3$.

There are at most $j$ choices for the missing input, so
\eqref{eq:Kmissing} bounds
$K_{\infty,j}-K_{\infty,j}^{D_N}$.  Finally, every nonzero folded output in
\eqref{eq:Kfold} equals $(p_0,p+2n\mathfrak R_N)$ and has spatial modulus at
least $\mathfrak R_N$; its lattice product is envelope-dominated, so the proof
of \eqref{eq:Khigh} bounds the finite sum of all such aliases.  Adding this
alias estimate and \eqref{eq:Kmissing} to \eqref{eq:Kexactdiff} proves
\eqref{eq:Kfolddiff}--\eqref{eq:Kcrossdiff}.
\end{proof}

\begin{proposition}[shifted-interaction comparison]\label{prop:shifted}
Let $\Delta\U_N:=\U_N(\cdot+h^{(N)})-\U_N(\cdot)$ denote the V1
Wick-binomial shift of the interaction,
\[
 \Delta\U_N=\lambda\,\eps_N\sum_x\int_0^t
 \Bigl[4\,h^{(N)}{:}\eta_N^3{:}_{c}+6\,(h^{(N)})^2{:}\eta_N^2{:}_{c}
 +4\,(h^{(N)})^3\eta_N+(h^{(N)})^4\Bigr]ds,
 \qquad c=c_{N,L},
\]
and $\Delta\U$ its continuum analogue with $h^{(\infty)}$ and the matched
convention (each chaos existing by the truncation argument of
Proposition~\ref{prop:Yexists} with deterministic bounded coefficients).
Then, for $N=M+r$, $\eps_N\le\tfrac12$,
\[
 \|\Delta\U_N-\Delta\U\|_{L^2(\PP)}
 \ \le\ C\Bigl[2^{-r(1+\eta)/4}
 +\eps_N\bigl(1+\log(1/\eps_N)\bigr)\Bigr],
\]
with $C=C(m,\lambda,t_0,L,\eps_M,\bar P_M)\,(1+t)^{5/2}$.
\end{proposition}

\begin{proof}
Write $h_N=h^{(N)}$, $h_\infty=h^{(\infty)}$ and let
$I_N(F):=\eps_N\sum_x\int_0^tF(s,x)ds$, $I_\infty(F):=\iint F(s,x)dxds$.
The two Wick conversions needed here are exact polynomial identities,
\[
 {:}x^3{:}_{c}={:}x^3{:}_{v}+3(v-c)x,qquad
 {:}x^2{:}_{c}={:}x^2{:}_{v}+(v-c).
\]
Consequently, with $d_N=v_N-c_{N,L}$ and its continuum analogue,
\begin{align*}
 G_{N,3}&=4h_N,&
 G_{N,2}&=6h_N^2,&
 G_{N,1}&=4h_N^3+12d_Nh_N,&
 G_{N,0}&=h_N^4+6d_Nh_N^2,
\end{align*}
and the same formulas with $N$ replaced by $\infty$, satisfy the exact identities
\begin{equation}
 \begin{aligned}
 \Delta\U_N&=\lambda\left[I_N(G_{N,0})+
   \sum_{j=1}^3I_N(G_{N,j}{:}\eta_N^j{:}_{v_N})\right],\\
 \Delta\U&=\lambda\left[I_\infty(G_{\infty,0})+
   \sum_{j=1}^3I_\infty(G_{\infty,j}{:}\eta^j{:}_{v_\infty})\right].
 \end{aligned}
 \label{eq:shift-ensemble-exact}
\end{equation}
Thus no Wick-recombination remainder is left implicit.

Fix $j\in\{1,2,3\}$ and abbreviate the two stochastic integrals in
\eqref{eq:shift-ensemble-exact} by $T_{N,j}$ and $T_{\infty,j}$.  Isserlis gives
\begin{equation}
 \frac1{j!}\E[(T_{N,j}-T_{\infty,j})^2]
 =Q_{N,j}-2Q_{\times,j}^{\rm true}+Q_{\infty,j}.                  \label{eq:shift-Qtrue}
\end{equation}
The true mixed kernel is $(C_\times+R_N)^j$.  The asymmetric-band kernel
$C_\times$ need not be real on the edge column, so define the real surrogate
and the physical correction by
\[
 Q_{\times,j}:=\Re\ip{G_{N,j}}{\mathbf C_\times^jG_{\infty,j}},
 \qquad
 \mathcal P_{N,j}:=Q_{\times,j}^{\rm true}-Q_{\times,j}
 =\Re\ip{G_{N,j}}{[(\mathbf C_\times+\mathbf R_N)^j-
                   \mathbf C_\times^j]G_{\infty,j}}.
\]
The last equality uses that the actual covariance pairing
$Q_{\times,j}^{\rm true}$ is real.  Taking the real part can only decrease
the absolute value, so the kernel sup-norm estimate below applies without loss.
By Proposition~\ref{prop:laws}(ii),
$\|R_N\|_\infty\le\varrho_N:=3\cth(2t_0)\eps_N/(8L)$, and the envelope
sum gives $\|C_\times\|_\infty\le C(m,t_0)\log_N$.  Hence
\begin{align}
 \|(C_\times+R_N)^j-C_\times^j\|_\infty
 &\le j\varrho_N(\|C_\times\|_\infty+\varrho_N)^{j-1},\notag\\
 |\mathcal P_{N,j}|
 &\le C(m,t_0,L,\Sigma_M(p),\bar d)(1+t)^2
          \eps_N\log_N^{j-1}.                                   \label{eq:physical-edge-shift}
\end{align}
This is the complete physical-edge contribution; it is not included in any Fourier
alias term below.

We now decompose the surrogate expression exactly.  Write
\[
 G_{N,j}(s,x)=\frac1{2L}\sum_{p\in B_N}b_N(s,p)e^{ipx},\qquad
 G_{\infty,j}(s,x)=\frac1{2L}\sum_{q\in\Gamma_\infty}b_\infty(s,q)e^{iqx},
\]
and set $\widetilde b(p_0,p)=\int_0^tb(s,p)e^{-ip_0s}ds$.  Parseval is
\begin{equation}
 \|b\|_{2,\#}^2:=\frac1{2Lt}\sum_{p_0,p}|\widetilde b(p_0,p)|^2
 =\int_0^t\frac1{2L}\sum_p|b(s,p)|^2ds.                           \label{eq:bparseval}
\end{equation}
Directly expanding the kernel powers and performing the time, lattice-space, and
continuum-space integrals gives
\begin{align}
 Q_{N,j}&=\frac1{2Lt}\sum_{p_0,p\in B_N}
 K_{N,j}^{\rm fold}(p_0,p)|\widetilde b_N(p_0,p)|^2,               \label{eq:QNexact}\\
 Q_{\times,j}&=\Re\left\{\frac1{2Lt}\sum_{q_0,q\in\Gamma_\infty}
 K_{\times,j}(q_0,q)
 \overline{\widetilde b_N(q_0,[q]_N)}\widetilde b_\infty(q_0,q)\right\}, \label{eq:Qxexact}\\
 Q_{\infty,j}&=\frac1{2Lt}\sum_{q_0,q\in\Gamma_\infty}
 K_{\infty,j}(q_0,q)|\widetilde b_\infty(q_0,q)|^2.              \label{eq:Qinfexact}
\end{align}
In \eqref{eq:Qxexact} the summand is zero when $K_{\times,j}=0$.  The displayed
real part is essential: the negative Nyquist edge belongs to $D_N$ while the
positive edge does not, so conjugate exact outputs need not both occur in the
surrogate sum.

Extend $b_N$ by zero outside $B_N$ and put
$\delta b=b_N-b_\infty$ on $B_N$, $\delta b=-b_\infty$ outside.  Splitting
$q=p+2n\mathfrak R_N$ in \eqref{eq:Qxexact} and expanding the square in
$\langle\delta b,K_{\infty,j}\delta b\rangle$ gives the identity
\begin{align}
 Q_{N,j}-2Q_{\times,j}+Q_{\infty,j}
 &=\langle\delta b,K_{\infty,j}\delta b\rangle
 +\langle b_N,(K_{N,j}^{\rm fold}-K_{\infty,j})b_N\rangle_{B_N}\notag\\
 &\quad-2\Re\langle b_N,(K_{\times,j}-K_{\infty,j})
                    b_{\infty,B_N}\rangle_{B_N}
 -2\Re A_{\times,j},                                                \label{eq:shift-exact-partition}
\end{align}
where every bracket includes the factor $(2Lt)^{-1}$ and
\begin{equation}
 A_{\times,j}:=\frac1{2Lt}
 \sum_{\substack{p_0,\ p\in B_N,\ n\ne0}}
 K_{\times,j}(p_0,p+2n\mathfrak R_N)
 \overline{\widetilde b_N(p_0,p)}
 \widetilde b_\infty(p_0,p+2n\mathfrak R_N).                       \label{eq:Across}
\end{equation}
Only finitely many $n$ occur, because a $j$-fold output of $D_N$ lies in the
spatial interval $[-j\mathfrak R_N,j\mathfrak R_N]$.  Equations
\eqref{eq:shift-exact-partition}--\eqref{eq:Across} are an algebraic equality;
there is no unnamed tail term.

By \eqref{eq:Khigh} and Cauchy--Schwarz, each Fourier coefficient is repeated in
\eqref{eq:Across} at most $2j+1$ times, so
\begin{equation}
 |A_{\times,j}|\le C_jB_j''\eps_N^2\log_N^{j-1}
                  \|b_N\|_{2,\#}\|b_\infty\|_{2,\#}.              \label{eq:Across-bound}
\end{equation}
Using \eqref{eq:Ksup}--\eqref{eq:Kcrossdiff} in
\eqref{eq:shift-exact-partition} therefore yields
\begin{equation}
 |Q_{N,j}-2Q_{\times,j}+Q_{\infty,j}|
 \le B_j\|\delta b\|_{2,\#}^2
 +C\eps_N^2\log_N^2
   (\|b_N\|_{2,\#}+\|b_\infty\|_{2,\#})^2.                         \label{eq:shift-Qbound}
\end{equation}

We record the coefficient estimate rather than referring to a generic ``same tail
count''.  For a spatial sequence set
$\|a\|_{1,L}=(2L)^{-1}\sum_k|a_k|$,
$\|a\|_{2,L}^2=(2L)^{-1}\sum_k|a_k|^2$, and
$(a*_Lb)(p)=(2L)^{-1}\sum_ka_kb_{p-k}$.  Young's inequality reads
$\|a*_Lb\|_{2,L}\le\|a\|_{2,L}\|b\|_{1,L}$.  The coefficient of $h^n$ is
$a^{*_{L}n}$, folded modulo $2\mathfrak R_N$ on the lattice and exact in the
continuum, and
\begin{equation}
 a_N^{*_{L}n}-a_\infty^{*_{L}n}
 =\sum_{\ell=0}^{n-1}a_N^{*_{L}\ell}*_{L}(a_N-a_\infty)
     *_{L}a_\infty^{*_{L}(n-1-\ell)}.                               \label{eq:power-telescope}
\end{equation}
Lemma~\ref{lem:H}(i) gives uniform $\ell^1_L$ bounds.  Lemma~\ref{lem:H}(iii)
controls the middle factor in $\ell^2_L$.  A folded alias or an output outside
$B_N$ satisfies $|k_1+\cdots+k_n|\ge\mathfrak R_N$, hence some
$|k_i|\ge\mathfrak R_N/n$.  From the D4 envelope, for $1\le n\le4$,
\begin{equation}
 \frac1{2L}\sum_{|k|\ge\mathfrak R_N/n}|a_k^{(\infty)}(s)|^2
 +\frac1{2L}\sum_{\substack{k\in B_N\\|k|\ge\mathfrak R_N/n}}
       |a_k^{(N)}(s)|^2
 \le C(\eta,n)\bar P_M^2\,2^{-r(1+\eta)}.                         \label{eq:D4-power-tail}
\end{equation}
Indeed $|a_k|^2\le C\eps_M\bar P_M^2(1+\eps_M|k|)^{-2-\eta}$ and the
sum--integral comparison from the lower limit
$\eps_M\mathfrak R_N/n=\pi2^r/n$ gives
$\int_{\pi2^r/n}^\infty(1+u)^{-2-\eta}du
=(1+\pi2^r/n)^{-1-\eta}/(1+\eta)$.  Applying Young to
\eqref{eq:power-telescope}, and placing the forced factor of an alias in
\eqref{eq:D4-power-tail}, proves for every $n\le4$
\begin{equation}
 \|b_N^{(h_N^n)}-b_\infty^{(h_\infty^n)}\|_{2,\#}^2
 \le Ct\,2^{-r\theta_H},\qquad \theta_H=(1+\eta)/2,                \label{eq:power-coeff-diff}
\end{equation}
where the left side includes the common-band difference, all folded aliases, and
all missing continuum outputs.  The same calculation and
$|d_N-d_\infty|\le C\eps_N^2$ give, for the exact coefficients $G_{a,j}$,
\begin{equation}
 \|b_N\|_{2,\#}^2+\|b_\infty\|_{2,\#}^2\le C(1+t),\qquad
 \|\delta b\|_{2,\#}^2\le C(1+t)^3
       [2^{-r\theta_H}+\eps_N^4].                                  \label{eq:Gcoeff}
\end{equation}

Insert \eqref{eq:Gcoeff} into \eqref{eq:shift-Qbound}, add the physical-edge
term \eqref{eq:physical-edge-shift} through \eqref{eq:shift-Qtrue}, and use
\begin{equation}
 \eps_N\log_N^2=\eps_M2^{-r}
 [1+\log(1/\eps_M)+r\log2]^2
 \le C(\eps_M,\eta)2^{-r\theta_H},                                  \label{eq:edge-absorb}
\end{equation}
which holds because $1-\theta_H>0$.  We obtain
\begin{equation}
 \E[(T_{N,j}-T_{\infty,j})^2]
 \le C(1+t)^4[2^{-r\theta_H}+\eps_N^2\log_N^2],
 \qquad j=1,2,3.                                                     \label{eq:shift-chaos-final}
\end{equation}

For $j=0$, $I_N(G_{N,0})-I_\infty(G_{\infty,0})$ is the zero spacetime
coefficient of the two- and four-fold convolutions in $G_{N,0},G_{\infty,0}$.
Equations \eqref{eq:power-telescope}--\eqref{eq:D4-power-tail}, followed by
Cauchy--Schwarz in spacetime and $|d_N-d_\infty|\le C\eps_N^2$, give
\begin{equation}
 |I_N(G_{N,0})-I_\infty(G_{\infty,0})|
 \le C(1+t)^2[2^{-r\theta_H/2}+\eps_N^2].                            \label{eq:shift-zero}
\end{equation}
Finally sum the four terms in \eqref{eq:shift-ensemble-exact}, take square roots in
\eqref{eq:shift-chaos-final}, and use
$\theta_H/2=(1+\eta)/4$, $\log_N\ge1$, and $\eps_N\le1$.  Enlarging the
fixed-parameter constant and the harmless time power to $(1+t)^{5/2}$ gives exactly
the asserted estimate.
\end{proof}

\section{Uniform Nelson bounds (A5)}\label{sec:A5}

\begin{lemma}[pointwise lower bound]\label{lem:lower}
For every $N$ and every realization,
\begin{equation}
 \U_N\ \ge\ -6\lambda\,(2Lt)\,c_{N,L}^2
 \ \ge\ -c_1\,t\,\bigl(1+\log(1/\eps_N)\bigr)^2,
 \qquad c_1=c_1(\lambda,m,L),
 \label{eq:lower}
\end{equation}
and the same bound holds for the shifted interaction
$\U_N(\cdot+h)$ for any deterministic \emph{real} shift $h$ --- reality being
essential, the identity ${:}u^4{:}_c=(u^2-3c)^2-6c^2$ being one between real numbers.
The shifts used below, $h^{(N)}$ and $h^{(\infty)}$ of Definition~\ref{def:shift},
are real by part (a) of the remark following that definition.
\end{lemma}

\begin{proof}
$\U_N=\lambda\eps_N\sum_x\int{:}\eta_N^4{:}_{c_{N,L}}$ (Definition~\ref{def:U}
$=$ V1 form) and pointwise ${:}u^4{:}_c=(u^2-3c)^2-6c^2\ge-6c^2$; sum and
integrate ($\eps_N\cdot2r_N\cdot t=2Lt$). The Wick constant obeys
$c_{N,L}\le C(m,L)(1+\log(1/\eps_N))$ by the band sum
$\sum_{\Gamma_N}\gamma_N^{-1}\le C\log$. The bound is evaluated pointwise in
the field argument, hence is invariant under deterministic real shifts (for complex
$u$ the identity gives no lower bound at all: at $u=ia$ with $a$ real it reads
$a^4+6ca^2+3c^2$).
\end{proof}

\begin{theorem}[A5: uniform Nelson bounds]\label{thm:A5}
For every $p<\infty$ and $t\in[t_0,T]$:
\begin{equation}
 \sup_{N}\ \bigl\|e^{-\U_N}\bigr\|_{L^p(\PP)}\ <\ \infty ,
 \label{eq:A5}
\end{equation}
with a bound depending on $(p,\lambda,m,L,t_0,T,\eps_M)$ --- here and only here $p$ is
the Lebesgue exponent, the coarse momentum field of the symbol $\xi=\eps_Mq+ip_M$
being written $p_M$; the dependence on $\eps_M$ enters through the fourth entry of
\eqref{eq:scale}, which is imposed for every admissible $b$ and hence also in the
unshifted case, and it is harmless because $M$ is fixed once and for all before
$N\to\infty$ --- ; the same holds, with
$\U_N$ replaced by $\U_N(\cdot+h^{(N)})$, uniformly over coarse symbols
with $\Sigma_M(p)\le R$ and $\bar P_M\le R'$ (constants then also depend on
$R,R'$).
\end{theorem}

\begin{proof}
Write $b_N^*:=c_1t(1+\log(1/\eps_N))^2$, so $\PP[\U_N\le-b]=0$ for
$b>b_N^*$ by Lemma~\ref{lem:lower}. Fix $b\le b_N^*$ with
\begin{equation}
 L_b:=\bigl(b/2c_1t\bigr)^{1/2}
 \ \ge\ \max\bigl(2(1+\log2),\ 3,\ 4\log2/c_2,\ 1+\log(1/\eps_{M})\bigr)
 \label{eq:scale}
\end{equation}
(the four entries are exactly what the three steps below require: $2(1+\log2)$ gives
$\eps_{N_b}\le\tfrac12$, the standing hypothesis of Theorem~\ref{thm:A4a}; $3$ gives
$\eps_{N_b}\le e^{-1/2}$, where $\eps\mapsto\eps(1+\log(1/\eps))^{3/2}$ is increasing;
$4\log2/c_2$ gives $2q\ge2$ in the hypercontractive step; and
$1+\log(1/\eps_M)$ gives $N_b\ge M$, needed for $r_b:=N_b-M$ and $h^{(N_b)}$ in the
shifted case. Smaller $b$ are absorbed into the constant at the end, at the cost of a
factor $e^{pb_1}$ in the layer cake, $b_1$ the value of $b$ at equality above), and
choose the comparison scale $N_b\in[M,N]$ as the largest
dyadic scale with $1+\log(1/\eps_{N_b})\le L_b$. This exists and satisfies
$1+\log(1/\eps_{N_b})>L_b-\log2\ge L_b/2$: the value at $M$ is admissible
by \eqref{eq:scale}, the value at $N$ is inadmissible since $b\le b_N^*$
forces $1+\log(1/\eps_N)\ge\sqrt2\,L_b>L_b$, and each dyadic refinement
increases $1+\log(1/\eps)$ by exactly $\log2$, so the largest admissible
scale $N_b<N$ exists and the next scale's violation gives the lower bound.
In particular $\eps_{N_b}\le e^{1-L_b/2}$. Then $\U_{N_b}\ge-c_1tL_b^2=-b/2$ pointwise
(Lemma~\ref{lem:lower}), so
\[
 \{\U_N\le-b\}\ \subseteq\ \{\U_N-\U_{N_b}\le-b/2\}.
\]
By Theorem~\ref{thm:A4a} and the triangle inequality through $\U$, and
since $\eps\mapsto\eps(1+\log(1/\eps))^{3/2}$ is increasing on
$(0,e^{-1/2})$,
\[
 \rho_b:=\|\U_N-\U_{N_b}\|_{L^2}
 \le C\bigl[\eps_N\log_N^{3/2}+\eps_{N_b}\log_{N_b}^{3/2}\bigr]
 \le C_\rho\,e^{-c_2L_b}\,L_b^{3/2},
 \qquad c_2=\tfrac12 .
\]
$\U_N-\U_{N_b}$ is a sum of Wick chaoses of order $\le4$, so Nelson's
hypercontractive moment comparison applies in the sum-of-chaoses form of
Part I, Appendix C (proof of the $h$-uniform stability theorem,
\texttt{thm:stability}: write $Y=\Gamma(e^{-\tau})Z$, $Z=\sum_re^{r\tau}Y_r$,
use $\|\Gamma(e^{-\tau})\|_{2\to q}\le1$ for $q=1+e^{2\tau}$ and chaos
orthogonality; for order $\le4$ this yields
$\|Y\|_{L^{2q}}\le(2q-1)^{2}\|Y\|_{L^2}$), and with Chebyshev in $L^{2q}$,
\[
 \PP[\U_N\le-b]\ \le\
 \Bigl(\frac2b\Bigr)^{2q}(2q-1)^{4q}\rho_b^{2q}
 \le\Bigl[\frac2b\,(2q)^2\rho_b\Bigr]^{2q}.
\]
Choose $(2q)^2=e^{c_2L_b/2}$: the bracket is
$\le\frac{2C_\rho}b\,L_b^{3/2}\,e^{-c_2L_b/2}\le e^{-1}$ for
$b\ge b_1$, $b_1=b_1(\lambda,m,L,t_0,T,\eps_M)$, whence
\[
 \PP[\U_N\le-b]\ \le\ \exp\bigl(-e^{c_2L_b/4}\bigr)
 =\exp\Bigl(-\exp\bigl(\tfrac{c_2}4(b/2c_1t)^{1/2}\bigr)\Bigr),
\]
uniformly in $N$ (for $b>b_N^*$ the probability is zero, which is smaller).
Since $e^{c\sqrt b}$ grows faster than any polynomial in $b$, the tail
exponent eventually dominates $2pb$, so
\[
 \E\bigl[e^{-p\,\U_N}\bigr]
 =\int_{-\infty}^\infty pe^{pb}\,\PP[\U_N\le-b]\,db
 \le e^{pb_1}+\int_{b_1}^\infty pe^{pb}\,
 e^{-\exp(\frac{c_2}4(b/2c_1t)^{1/2})}\,db<\infty
\]
with an $N$-independent bound (layer-cake identity as in Part I, App.\ C;
the $b<b_1$ range contributes $e^{pb_1}$ since $\PP\le1$). For the shifted
version: Lemma~\ref{lem:lower} holds verbatim under any deterministic
shift, and $\U_N(\cdot+h^{(N)})-\U_{N_b}(\cdot+h^{(N_b)})
=(\U_N-\U_{N_b})+(\Delta\U_N-\Delta\U_{N_b})$, where the second bracket is
controlled by Proposition~\ref{prop:shifted} at the two scales: with
$r_b:=N_b-M$, $2^{-r_b(1+\eta)/4}=(\eps_{N_b}/\eps_M)^{(1+\eta)/4}
\le C(\eps_M)e^{-L_b(1+\eta)/8}$, so the total obeys the same bound with
$c_2=(1+\eta)/8$ and a constant uniform over $\Sigma_M(p)\le R$,
$\bar P_M\le R'$ (Lemma~\ref{lem:H}); the argument then runs verbatim with
this smaller $c_2$.
\end{proof}

\section{Acceptance and claim ledger}\label{sec:ledger}

V0 acceptance: Theorem~\ref{thm:dict} proves the constant identified
numerically ($0.0384461803\ldots$ to $3.6\times10^{-10}$); no other V0 item
was scheduled against V2. The A4/A5 outputs consumed by A6 are:
Theorem~\ref{thm:A4a}, Proposition~\ref{prop:shifted},
Theorem~\ref{thm:A5}, and Lemma~\ref{lem:recomb}/\ref{lem:H}; the
convention lock for the gate's right-hand side is
Corollary~\ref{cor:lock}.

\begin{longtable}{P{5.6cm}P{2.4cm}P{5.8cm}}
\toprule
claim & status & ground \\
\midrule
\endhead
Same-noise coupling realizes the V1 lattice ensemble & proved &
Prop.~\ref{prop:laws} (Mittag--Leffler sum \eqref{eq:ML}) \\
Recombination; $d_N(t)$ finite, $\to d_\infty(t)$, $O(\eps_N^2)$ & proved &
Lemma~\ref{lem:recomb} \\
Existence of the continuum Wick integrals $Y_2,Y_4$ & proved &
Prop.~\ref{prop:Yexists} (monotone mode sums) \\
Envelope $\hat c\le\frac{\pi^2/4}{|\kappa|^2+m^2}$; convolution decay;
summability & proved & Lemma~\ref{lem:envelope} \\
Alias sums $\le C(m,t_0)tL\,\eps_N^2\log_N^3$ & proved &
Lemma~\ref{lem:alias} \\
Symmetric-edge correction: one-, two-, and four-edge sectors retained;
exactly three impossible; $\Delta_{4e}\le
\frac{9t}{256L^2}\cth(2t_0)^3\eps_N^5$ & proved &
Proposition~\ref{prop:laws}, equations \eqref{eq:four-edge}--\eqref{eq:Eedge-bound} \\
$\|\U_N-\U\|_{L^2}\le C(1+t)(1+L)\,\eps_N\log_N^{3/2}$ & proved &
Theorem~\ref{thm:A4a} (second chaos: exact-square structure) \\
Dictionary $\delta_\infty=\frac1{2\pi}\log\frac4\pi$ with rate
$O(\eps_N)$, explicit constant & proved & Theorem~\ref{thm:dict}; V0
confirmation \\
Convention lock: gate polynomial is plain $\varphi^4$ (no counterterm;
the two tadpole offsets cancel) & proved & Corollary~\ref{cor:lock},
Remark~\ref{rem:deltarole} \\
Closed form $H(s)=\frac{\rho}2\,
\frac{\sh(\omega(s-t/2))}{\sh(\omega t/2)}$; $|H|\le|\rho|/2$;
$\omega$-derivative bound & proved & Lemma~\ref{lem:closed} \\
Uniform coarse-data shift bounds ($\Sigma_M(p)$); symbol convergence at
rate $2^{-r(1+\eta)/2}$ & proved & Lemma~\ref{lem:H} \\
Folded lattice, exact-output mixed, and continuum power symbols: uniform
bounds, fundamental-band differences, and high-output tails & proved &
Lemma~\ref{lem:symbolsup}, equations \eqref{eq:Kfolddiff}--\eqref{eq:Khigh} \\
Shifted-interaction $L^2$ comparison,
$O(2^{-r(1+\eta)/4}+\eps_N\log_N)$; exact four-term quadratic partition
for the real asymmetric-band surrogate, with explicit $A_{\times,j}$ & proved & Prop.~\ref{prop:shifted},
equations \eqref{eq:shift-exact-partition}--\eqref{eq:Across} \\
Uniform Nelson bounds, plain and shifted & proved & Theorem~\ref{thm:A5} \\
\bottomrule
\end{longtable}

\medskip
\noindent\emph{External imports used (verbatim, with source locations):}
\begin{itemize}
\item free-rate document (file \texttt{lamport\_part\_ii\_free\_fixed\_%
coarse\_rate}): Lemma 2.1, eq.\ (2.4) ($D4$ envelope
$|P_r(\ell)|^2\le C_D(1+|\ell|)^{-2-\eta}$, $\eta=2-\log_23$, proved for
$|\ell|\le\pi2^r$ and, for $P_\infty$, on all of $\R$); Lemma 3.1, eq.\ (3.1)
($0\le\gamma-\gamma_N\le\eps_N^2k^4/(24m)$, under $|\eps_Nk|\le\pi$); and the display
(3.5) \emph{inside the proof} of that lemma (its step LS 1.1),
$0\le k^2-4\eps_N^{-2}\sin^2(\eps_Nk/2)\le\eps_N^2k^4/12$, likewise valid under
$|\eps_Nk|\le\pi$ --- which holds throughout since $k\in\Gamma_N$ --- consumed in
Lemma~\ref{lem:symdiff};
\item model sheet (file \texttt{lamport\_part\_ii\_model\_sheet}): the conventions of
\S1--\S2 and, in Lemma~\ref{lem:H}, the mask coefficients (2.1), from which
$L_{m_0}\le\tfrac1{\sqrt2}\sum_n n|h_n|=(3+\sqrt3)/4$;
\item V1 (file \texttt{v1\_thermal\_representations}): Lemma 5.2
(\texttt{lem:interp}, interpolant characterization and uniqueness),
Theorem 5.4 (\texttt{thm:meanshift}, Wick-binomial shifted form), Remark
5.5 (\texttt{rem:ensemble}, the periodic ensemble matched in
Proposition~\ref{prop:laws});
\item Part I (file \texttt{part\_i/main.tex}), Appendix C, proof of the
$h$-uniform stability theorem (labels \texttt{app:huniform},
\texttt{thm:stability}): the $\Gamma(e^{-\tau})$ hypercontractive moment
comparison for sums of Wick chaoses, and the layer-cake assembly of
$\E[e^{-p\,\U}]$ from the doubly-exponential tail.
\end{itemize}
The $D4$ mask Lipschitz constant
$L_{m_0}\le(3+\sqrt3)/4$ is derived inline in Lemma~\ref{lem:H} from the
model-sheet mask coefficients. No other external facts are consumed.

\noindent\emph{Phase V2 gate:} A4(a), A4(b), A4(c), and A5 are closed;
constants explicit; the $\delta_\infty$ value matches V0. Remaining phases:
V3 (Statement A1 with the (A1-iv) FK clause) and V4 (A6--A8 assembly and the
adversarial pass).

\end{document}
